% Modernised reading — fonds Alexandre Grothendieck, shelfmark 78, pages 1-129
% (the whole folder).
%
% This is an INTERPRETATION, not a transcription. It re-reads the pages in
% today's notation and vocabulary, states the hypotheses the manuscript leaves
% implicit, and drops the critical apparatus so that the mathematics reads
% straight through. Everything the brackets and underlines carried moves into
% a footnote. It is held to being mathematically correct as it stands; where
% the manuscript is loose or mistaken, the true statement is what appears here
% and the footnote says what the page has. In any dispute of substance the
% transcription governs: whoever cites Grothendieck cites batch-01.fr.tex to
% batch-07.fr.tex (pp. 1-20, 21-40, 41-60, 61-80, 81-100, 101-120, 121-129).
%
% Pass: Opus 5.5 (claude-opus-5-5), 2026-10-10 — first pass, unchecked against the pages by a human.
% The folder was read whole, all seven batches together; this is one document
% for the shelfmark. It was made on Opus 5.5 and has not been measured against
% the reference reading of folder 115 (made on Opus 5). The computations the
% reading asserts beyond the page (the matrices of sigma_0, sigma_1 and the
% condition sigma_0^2 = id, the centre, the invariant conic and its
% discriminants, the vertices s_n = (F_n F_{n+1}, F_{n-1} F_n), the identity
% X_n - U X_{n-1} + X_{n-2} = 1 and the Bezout pair it gives, the values
% F_n(±2), the table of Psi_n, the strip and parabola cases, the affine
% extension of i -> ri for alpha = ±2) were checked by hand while writing,
% not by machine. The literature named in footnotes (Lehmer on the minimal
% polynomial of 2cos(2pi/n), affinely regular polygons in finite planes,
% Chebyshev and Dickson polynomials, Giraud's gerbes) is cited from memory
% and was not checked against the sources.
%
% Scope. Pages not transcribed and not read here: 49 (a cover, « Cours /
% Polygones réguliers / II », which heads a section), 53 (blank letterhead),
% 93 (a typed letter to Grothendieck dated Paris, 10 November 1977; its date
% only is used, as the transcription's note does), 108 (a pink cover,
% « Polynômes cyclotomiques »). Page 110 is an unrelated typescript; only the
% two lines of his calculation at its foot are used. Page 114's lower half
% carries the start of an unrelated letter, not transcribed.
%
% Spine. One object, the regular polygon in an affine plane over an arbitrary
% field (then ring), classified by one number alpha = zeta + zeta^{-1}:
%   pp. 1-48    « Chap I », first draft: definitions (regular, unimodular,
%               proper), the real examples, pinnings, dihedral
%               representations (Props 1-3), dropping faithfulness (Prop 4),
%               infinite polygons (Prop 5), « étoilage » (Prop 6),
%               quasi-polygonal structures and regular systems of vertices
%               (Prop 7), a programme of descent over Z[1/n];
%   pp. 49-80   « Cours Polygones réguliers II », second draft: combinatorial
%               polygons as D_infinity-sets, orientations, the case k
%               algebraically closed with p not dividing n, invariant forms
%               and the circumscribed conic, the programme in « universal
%               affine coordinates », the Lemma and Proposition of p. 80;
%   pp. 81-100  the universal coordinates: (alpha, beta), proper and improper
%               realisations, fixed points and centre, eigenvalues, the three
%               cases alpha != ±2, alpha = -2, alpha = 2, the invariant
%               quadratic form f_alpha and its conic;
%   pp. 101-107 faithfulness, in two versions (p. 104 marked « ancienne
%               version »);
%   pp. 108-120 « Polynômes cyclotomiques »: F_n, the theorem
%               u_alpha^n = id <=> F_n(alpha) = 0, the semi-cyclotomic Psi_n,
%               the « modular rings », Phi_n mod p and the order of a polygon
%               in characteristic p, a table of alpha;
%   pp. 121-129 loose sheets: a struck tensor computation, tables of Psi_n
%               and Phi_n with their reductions mod p, the rings
%               Z[zeta], Z[zeta + zeta^{-1}] and their ramification at p, the
%               involution T -> T^{-1}, and a last page on ramified coverings
%               that breaks off.
%
% Notation fixed for the whole document.
%  - k a field (a ring where said), E an affine plane over k, V its direction,
%    Aff(E) the affine group; Dr(E) the set of lines of E.
%  - D_n the dihedral group of order 2n (2 <= n <= infinity), presented by two
%    involutions sigma_0, sigma_1, with u = sigma_0 sigma_1 the rotation of
%    one step. sigma_1 is the reflection that fixes the base vertex s_0,
%    sigma_0 the one that fixes the base side a_0 = dr(s_0, s_1). This is
%    the convention of pp. 77-120, where all the computations are. Chap I
%    writes sigma for sigma_1 and sigma' for sigma_0 (p. 17: u = sigma'sigma);
%    pp. 55-57 take u = sigma_1 sigma_0, the inverse rotation (footnoted);
%    p. 29 calls the vertex-fixing involution sigma_0 (footnoted).
%  - s_i = u^i s_0, a_i = dr(s_i, s_{i+1}); Pi_n the standard combinatorial
%    n-gon (the page's Pi_0(n) in Chap I, Pi_n from p. 57).
%  - alpha the « invariant numérique », zeta, zeta^{-1} the eigenvalues of
%    the rotation, alpha = zeta + zeta^{-1}.
%  - omega^{(r)} = {u^r, u^{-r}} the « étoilage » of a polygonal structure
%    (p. 13's own notation); the pages 33-39 write Psi_r(omega). Q_n =
%    (Z/nZ)^*/{±1} the group acting (the page's Psi_n, p. 33). Both renamed
%    because Psi_n is, from p. 114, the semi-cyclotomic polynomial.
%  - F_n in Z[U] the polynomials of p. 109 (product over the divisors d >= 3
%    of n of the Psi_d); Psi_n the minimal polynomial of zeta + zeta^{-1},
%    zeta a primitive n-th root of unity (p. 114). Pages 122-124 call the
%    latter F_n; it is written Psi_n here throughout, with a footnote.
%  - S_n in Z[U] with S_n(T + T^{-1}) = T^n + T^{-n} (n >= 1), S_0 = 1
%    (p. 122).
%
% Substantive departures from the pages, each footnoted where it occurs:
%  p. 3    unimodularity is « automatically satisfied » except, by pp. 84-85,
%          for the improper 2p-gons in characteristic p, which are regular
%          and not unimodular; the illegible word there is not restored.
%  p. 7    the classes of real regular polygons are T/{±1} minus the classes
%          of 1 AND -1, not of 1 alone.
%  pp. 22-23 Prop 4: in (b) the divisor d must be proper; in (c) the kernel
%          is generated by u^{n'}, n' = card S >= 3, n' | n (the page mixes
%          d and n' = n/d).
%  p. 27   the check « 3 <= d <= [n/3] » misses the doubled polygon: the
%          first non-faithful case is n = 6 (a triangle traversed twice), not
%          n = 9; the condition is vacuous exactly for n an odd prime or 4.
%  p. 33   Prop 6 a): for n infinite, r = ±1 also works.
%  p. 41   Aff±(S) has index phi(n)/2 in Aff(S), not 2.
%  pp. 44-45 Prop 7's condition 2) must keep the struck clause (sigma inverts
%          Z) and sigma != 1, or it is empty; unimodularity is not repeated.
%  p. 57   S_n and A_n are interchanged with respect to p. 54's S = R/sigma_1.
%  p. 63   the alternative zeta zeta' = -1 ({1, -1}) is excluded by n >= 3.
%  p. 66   the basis of invariants needs (xy)^a alone (i = 0) as well.
%  p. 97   « alpha != 0, 1 »: the triangle is alpha = -1; uniqueness of the
%          circumscribed conic holds for n >= 5.
%  p. 111  the third column of the matrix of u^n is (F_n F_{n+1},
%          F_{n-1} F_n), as on the struck first version.
%  p. 112  the conclusion is F_n(alpha) = 0 (page: = 1).
%  p. 113  A_n = F_{n-2} - U F_n (the page's F_{n-1} is doubtful and wrong).
%  p. 116  alpha = -2 is the invariant of the 2p-gon, not of a 2-gon.
%  pp. 118-119 Corollary: for n' = 1 and p = 2 the order is 4, not p.
%  p. 122  the recurrence for S_n is right as written, because S_0 = 1
%          (the transcription's note asks for a 1/2 that S_0 = 1 absorbs).
%  pp. 123-124 (p, n') = 1, not (n, n') = 1; F_p = Psi_p = (T - 2)^{p'}
%          mod p (page: T - 1), Psi_{2p} = (T + 2)^{p'}; « n != 4 » should
%          read « n != 2 ».
%
% Modern names given and footnoted as ours: affinely regular polygon, flag,
% torsor, gerbe, group scheme mu_n semidirect Z/2, affine line over Z/nZ,
% Chebyshev and Dickson polynomials, minimal polynomial of 2cos(2pi/n),
% maximal real subfield, total ramification, casus irreducibilis, Minkowski's
% theorem (p. 129). The pages name a person once (p. 70, read with doubt);
% nothing on them cites the literature.
%
% Announced and never established: the « question longue » of p. 11 (taken
% up and answered on pp. 81-111); proof of Prop 4 (b) (p. 23); the meaning of
% the degenerate étoilé representations (p. 39); the case Aff(S) of p. 44;
% condition 3) of Prop 7; the « Construction » of p. 48 (SRnS over Z or
% Z[1/2]); the Exercise of p. 61; the theory of (étale) morphisms and
% quotients Pi_(n) of combinatorial polygons (pp. 106-107); the question of
% p. 127 (answered on p. 128); the connectedness question of p. 129.
\documentclass[11pt,a4paper]{article}
\input{../preamble/grothendieck}

\folder{78}
\pages{1}{129}
\dating{[à partir de 1977]}
\watermark{Édition de démonstration\\interprétation personnelle de l'œuvre}
\foldertitle{[Polygones réguliers et polynômes cyclotomiques] : notes manuscrites (s.d.) — lecture modernisée du dossier entier}

\begin{document}

\section*{Résumé}

\begin{resume}
Ces feuillets se demandent ce qu'est un polygone régulier quand on retire au
plan tout ce qui permet d'ordinaire de le reconnaître : les longueurs, les
angles, et même les nombres réels. Il reste un plan « affine » --- des points,
des droites, le parallélisme --- sur un corps de nombres quelconque, y compris
un corps fini où $1 + 1 + 1 = 0$. Qu'est-ce qu'un pentagone régulier dans un
tel plan, et combien y en a-t-il ?

La réponse commence par un changement de définition. Au lieu de « côtés égaux
et angles égaux », qui n'a plus de sens, Grothendieck demande que toute
symétrie du dessin combinatoire --- tourner d'un cran, retourner --- soit
réalisée par une transformation du plan. C'est la définition par les
symétries, et elle est plus généreuse qu'on ne croit : dans le plan affine,
tout triangle est régulier, tout parallélogramme aussi, puisqu'une
transformation affine envoie n'importe quel triangle sur n'importe quel autre.
Un polygone régulier devient ainsi la donnée d'une représentation du groupe
diédral --- deux symétries $\sigma_0$, $\sigma_1$ dont le produit est la
rotation d'un cran --- et d'un sommet de départ.

L'idée qui arrive ensuite est qu'un seul nombre suffit à tout décrire. Dans le
plan réel, un $n$-gone régulier (étoilé ou non) est déterminé, à transformation
affine près, par $\alpha = 2\cos(2\pi k/n)$, et la règle qui fabrique les
sommets de proche en proche est la même pour tous : quatre sommets consécutifs
vérifient $s_{i+3} - s_i = (1 + \alpha)(s_{i+2} - s_{i+1})$. Grothendieck prend
ce nombre comme seule donnée, sur un corps quelconque : chaque valeur de
$\alpha$ définit un « polygone » infini, et la question devient de savoir
quand il se referme après $n$ pas. La réponse est une équation polynomiale à
coefficients entiers en $\alpha$ --- au fond, le polynôme minimal de
$2\cos(2\pi/n)$, cousin des polynômes cyclotomiques. Les $n$-gones réguliers
d'un corps correspondent ainsi aux racines de ce polynôme dans ce corps, et
l'anneau qu'on obtient en adjoignant une racine aux entiers est l'anneau « sur
lequel s'écrit » le $n$-gone régulier le plus général.

Le plus beau est ce qui se passe en caractéristique $p$. Les deux valeurs
$\alpha = 2$ et $\alpha = -2$, qu'on exclut d'habitude parce qu'elles
correspondent à une « rotation d'angle nul », y donnent de vrais polygones
réguliers : pour $\alpha = 2$, un polygone à $p$ sommets posés sur une
parabole, où la rotation devient une translation le long de la parabole ; pour
$\alpha = -2$, un polygone à $2p$ sommets qui zigzague entre deux droites
parallèles. Le cercle circonscrit devient une parabole, ou une paire de
droites ; la « métrique » invariante dégénère. Ces cas sont, dit-il, « les
seuls cas vraiment nouveaux », et ceux qu'on perdrait si l'on exigeait dès le
départ une métrique non dégénérée.

Le dossier se lit aussi pour sa manière. Grothendieck refuse de « traîner tout
au long » une condition (la fidélité) qu'on n'examine utilement qu'à certains
moments ; il préfère étudier d'un coup tous les polygones, finis et infinis,
puis demander lesquels se referment ; il garde les cas qui paraissent
« exotiques » parce qu'ils « reviendraient indirectement » de toute façon, et
qu'ils permettent de suivre la variation de l'objet avec son paramètre. Les
deux rédactions successives --- un premier chapitre, puis des notes de cours
--- montrent la même matière reprise depuis les définitions, plus simple la
seconde fois.

Les noms à chercher ensuite : polygones affinement réguliers (notamment dans
les plans affines finis), représentations du groupe diédral, polynôme minimal
de $2\cos(2\pi/n)$ et sous-corps réel maximal d'un corps cyclotomique,
polynômes de Tchebychev et de Dickson, ramification de $p$ dans
$\mathbf{Z}[\zeta_p]$, torseurs et gerbes, schéma en groupes $\mu_n \rtimes
\mathbf{Z}/2$.

\keywords{affinely regular polygon, dihedral group, representation of the infinite dihedral group, flag, torsor, regular polygon in positive characteristic, parabola, invariant quadratic form, circumscribed conic, cyclotomic polynomial, minimal polynomial of 2cos(2pi/n), maximal real subfield, Chebyshev polynomials, Dickson polynomials, ring of integers of a cyclotomic field, total ramification, affine line over Z/nZ, dihedral group scheme, gerbe, descent, casus irreducibilis}
\end{resume}

\pagerange{1}{129}
\section*{Le fil du dossier, et les conventions}

La chemise porte au crayon « Polygones réguliers / Chap I » ; l'inventaire
titre « [Polygones réguliers et polynômes cyclotomiques] », d'après une
seconde chemise, rose, insérée à la page 108. Une lettre dactylographiée
reçue, datée de Paris le 10 novembre 1977 (page 93, non transcrite), donne un
terme a quo compatible avec la datation de l'inventaire. Les pages 49 et 108
sont des chemises, la page 53 une feuille blanche à en-tête, les pages 110
et 114 portent au revers ou au pied des textes sans rapport.

Le dossier se rattache au cours de DEA de 1977--1978, « géométrie arithmétique
des polygones réguliers », dont l'annonce est au dossier 75 (page 32) : la
notion de polygone régulier sur un corps, puis sur un anneau de base,
« l'anneau sur lequel s'écrit le $n$-gone régulier le plus général », les formes
imprévues que prennent en caractéristique positive la réflexion et le centre.
C'est exactement la matière qu'on trouve ici. La page 49 est la chemise d'une
deuxième leçon, « Cours / Polygones réguliers / II ».

\subsection*{Les stations}

\begin{itemize}
  \item \emph{Pages 1 à 48 : « Chap. I », première rédaction.} Définitions
    (polygone régulier, unimodulaire, propre ; pages 2 à 5) ; les exemples
    réels $P_\zeta$ (pages 6 à 8, et une feuille de calculs, page 9) ;
    repères et épinglages (pages 10 et 11) ; les représentations du groupe
    diédral, propositions 1 à 3 (pages 12 à 21) ; abandon de la condition de
    fidélité, proposition 4 (pages 22 à 26) ; les $\infty$-polygones,
    proposition 5 (pages 26 à 33) ; l'étoilage, proposition 6 (pages 33 à
    39) ; structures quasi-polygonales et systèmes réguliers de sommets,
    proposition 7, et un programme de descente (pages 39 à 48).
  \item \emph{Pages 49 à 80 : « Cours Polygones réguliers II », seconde
    rédaction.} Plan et rappels (pages 50 à 52) ; polygones combinatoires,
    contours et drapeaux, comme ensembles sous le groupe diédral infini
    (pages 54 à 61) ; le cas d'un corps algébriquement clos de
    caractéristique première à $n$, formes invariantes, conique circonscrite
    (pages 62 à 73) ; reformulation du problème pour le polygone infini
    $\Pi_\infty$, lemme et proposition (pages 74 à 80).
  \item \emph{Pages 81 à 100 : les coordonnées universelles.} Les paramètres
    $(\alpha, \beta)$ et la condition $\sigma_0^2 = 1$ ; cas propre et cas
    impropre ; points fixes et centre ; valeurs propres ; les trois cas
    $\alpha \neq \pm 2$, $\alpha = -2$, $\alpha = 2$ ; la forme quadratique
    invariante $f_\alpha$ et sa conique ; la parabole.
  \item \emph{Pages 101 à 107 : la fidélité}, en deux versions, la seconde
    (page 104) marquée « ancienne version ».
  \item \emph{Pages 108 à 120 : « Polynômes cyclotomiques ».} Les polynômes
    $F_n$ ; le théorème $u_\alpha^n = 1 \iff F_n(\alpha) = 0$ ; le polynôme
    semi-cyclotomique $\Psi_n$ et les « anneaux modulaires » ; réduction de
    $\Phi_n$ modulo $p$ et ordre d'un polygone en caractéristique $p$ ; table
    des $\alpha$.
  \item \emph{Pages 121 à 129 : feuilles détachées.} Un calcul tensoriel barré
    (page 121) ; tables de $\Psi_n$ et $\Phi_n$, réductions modulo $p$ (pages
    122 à 124) ; les anneaux $\mathbf{Z}[\zeta]$ et $\mathbf{Z}[\zeta +
    \zeta^{-1}]$, ramification en $p$ (pages 125 à 127) ; l'involution $T
    \mapsto T^{-1}$ (page 128) ; revêtements ramifiés (page 129), où le
    dossier s'arrête.
\end{itemize}

Les deux premières stations sont deux \emph{moutures} de la même théorie,
non une suite : le cours reprend depuis les définitions ce que le chapitre
avait construit, avec un autre langage (le polygone combinatoire y devient un
ensemble de drapeaux sous le groupe diédral infini) et un programme plus net
(tout déduire du polygone infini). On les garde distinctes ; on renvoie de
l'une à l'autre là où elles se recouvrent. De même pour les deux versions de
la fidélité (pages 101 à 103 et 104 à 107) et pour les deux présentations des
polynômes (pages 109 à 120 et 122 à 124), qui n'appellent pas les mêmes
objets du même nom.

\subsection*{Conventions}

$k$ est un corps (un anneau quand on le dit), $E$ un plan affine sur $k$, $V$
son espace des translations, $\mathrm{Aff}(E)$ son groupe affine,
$\mathrm{Dr}(E)$ l'ensemble de ses droites ; pour un automorphisme affine $g$,
$g_V$ est sa partie linéaire. $\mathbf{D}_n$ ($2 \leq n \leq \infty$) est le
groupe diédral d'ordre $2n$,
\[
  \mathbf{D}_n = \langle \sigma_0, \sigma_1 \mid \sigma_0^2 = \sigma_1^2 =
  (\sigma_0\sigma_1)^n = 1 \rangle , \qquad
  \mathbf{D}_\infty = \mathbf{Z}/2\mathbf{Z} * \mathbf{Z}/2\mathbf{Z} ,
\]
et $u = \sigma_0 \sigma_1$ est la rotation d'un cran ; $\mathbf{D}_n^{+} =
\langle u \rangle$ le sous-groupe des rotations, d'indice $2$. Dans un
polygone épinglé, $\sigma_1$ est la symétrie qui fixe le sommet de base
$s_0$, $\sigma_0$ celle qui fixe le côté de base $a_0$ ; $s_i = u^i s_0$ et
$a_i = \mathrm{dr}(s_i, s_{i+1})$\footnote{C'est la convention des pages 77
à 120, où sont tous les calculs. Le chapitre I écrit $\sigma$ pour
$\sigma_1$ et $\sigma'$ pour $\sigma_0$ (page 17 : « $\sigma' = u\sigma$ tel
que $u = \sigma'\sigma$ »), les pages 55 à 57 prennent $u = \sigma_1\sigma_0$,
la rotation inverse, et la page 29 appelle $\sigma_0$ l'involution qui fixe un
sommet. Rien n'en dépend : on le signale là où cela se voit.}. $\Pi_n$ est le
$n$-gone combinatoire type.

L'\emph{invariant numérique} $\alpha$ est la trace de la rotation $u_V$ ; ses
valeurs propres sont $\zeta$, $\zeta^{-1}$, $\alpha = \zeta + \zeta^{-1}$.
L'\emph{étoilage} d'une structure polygonale $\omega = \{u, u^{-1}\}$ est
$\omega^{(r)} = \{u^r, u^{-r}\}$, et $Q_n = (\mathbf{Z}/n\mathbf{Z})^{*}/\{\pm
1\}$ le groupe qui agit par étoilage\footnote{La page 13 écrit elle-même
$\omega^{(r)}$ ; les pages 33 à 39 écrivent $\Psi_r(\omega)$ et notent
$\Psi_n$ le groupe $Q_n$. On renomme les deux, parce qu'à partir de la page 114
$\Psi_n$ est le polynôme semi-cyclotomique.}.

Deux familles de polynômes de $\mathbf{Z}[U]$, $U = T + T^{-1}$, se croisent
dans le dossier, et il faut les distinguer d'emblée :
\begin{itemize}
  \item $F_n$ (page 109), dont les racines sont les $\zeta + \zeta^{-1}$ pour
    \emph{toutes} les racines $\zeta \neq \pm 1$ de $\zeta^n = 1$ ; c'est le
    polynôme qui dit quand le polygone d'invariant $\alpha$ se referme après
    $n$ pas ;
  \item $\Psi_n$ (page 114), le polynôme minimal de $\zeta + \zeta^{-1}$ pour
    $\zeta$ racine \emph{primitive} $n$-ième ; il dit quand le polygone a
    exactement $n$ côtés.
\end{itemize}
On a $F_n = \prod_{d \mid n,\ d \geq 3} \Psi_d$\footnote{Cette factorisation
est de nous ; elle se lit sur les racines, et sur la table de la page 109
($F_6 = U^2 - 1 = \Psi_3 \Psi_6$). Les pages 122 à 124 appellent $F_n$ ce
qui est ici $\Psi_n$ ; on écrit $\Psi_n$ partout.}. Enfin $S_n \in
\mathbf{Z}[U]$ est défini par $S_n(T + T^{-1}) = T^n + T^{-n}$ pour $n \geq 1$
et $S_0 = 1$.

Le mot « propre » a deux sens dans le dossier : pour un polygone (page 3),
l'incidence géométrique reflète l'incidence combinatoire --- ce que le cours
appelle « fidèle » (page 50) ; pour une réalisation de $\Pi_\infty$ (page 84),
la réalisation n'est pas le cas « impropre » $\beta = -2$. On précise chaque
fois.

\subsection*{Ce que le dossier annonce sans l'établir}

La démonstration de la partie (b) de la proposition 4 (page 23) ; le sens des
représentations étoilées dégénérées (page 39) ; la détermination de la
structure quasi-polygonale par $\mathrm{Aff}(S)$ (page 44, « faire dém. ») ;
la condition 3) de la proposition 7 ; la « Construction » d'un système
régulier de sommets sur $\mathbf{Z}$ ou $\mathbf{Z}[\frac12]$ (page 48) ;
l'exercice de la page 61 ; la théorie des morphismes (étales) de polygones
combinatoires et des quotients $\Pi_{(n)}$ dont la page 106 a besoin ; la
question de connexité de la page 129. La « question longue » de la page 11 ---
écrire explicitement le plan et les sommets en fonction de $n$ et d'un module
$\alpha$ --- est, elle, menée à bien aux pages 81 à 111.

\pagerange{1}{48}
\section*{I. Le chapitre I : polygones réguliers affines (pages 1 à 48)}

\pagerange{1}{5}
\subsection*{Définitions (pages 1 à 5)}

Soit $E$ un plan affine sur un corps $k$. Un \emph{polygone} (« déployé »)
dans $E$ est une partie finie $S$ de $E$, l'ensemble des sommets, munie d'une
\emph{structure polygonale combinatoire} : une partie $\omega = \{u,
u^{-1}\}$ du groupe symétrique $\mathfrak{S}_S$ formée de deux permutations
circulaires inverses l'une de l'autre, ce qui force $n = \operatorname{card} S
\geq 3$. Les droites qui joignent deux sommets consécutifs sont les
\emph{côtés} ; $n$ est aussi le nombre de côtés combinatoires, et l'on parle de
$n$-gone\footnote{La paire $\{u, u^{-1}\}$ est la « bi-circulation » du
dossier 81 (pages 41 et 42) et la troisième description d'un polygone du
dossier 89 (pages 19 et 20). Ici le mot « orientation » désigne l'un des deux
éléments de $\omega$.}.

\textbf{Polygone régulier.} $P = (E, S, \omega)$ est \emph{régulier} si tout
automorphisme combinatoire de $(S, \omega)$ --- tout élément du normalisateur
$\Gamma$ de $\omega$ dans $\mathfrak{S}_S$, groupe diédral d'ordre $2n$ --- se
prolonge \emph{de façon unique} en un automorphisme affine de $E$. Il suffit
de l'exiger pour un système de générateurs : $u$ et la symétrie $\sigma_s$
fixant un sommet $s$, ou $\sigma_s$ et la symétrie $\sigma_a$ fixant un côté
$a$ ; si $s$ et $a$ sont incidents, deux quelconques des trois éléments
$\sigma_s$, $\sigma_a$, $u = \sigma_a \sigma_s$ engendrent $\Gamma$.
L'unicité du prolongement équivaut à ce que $S$ engendre affinement $E$, ou
encore à ce que trois sommets consécutifs $s_0, s_1, s_2$ ne soient pas
alignés : s'ils l'étaient, la régularité donnerait $u(d_0) = d_0$ pour la
droite $d_0 = \mathrm{dr}(s_0, s_1)$, et $u$ étant transitif sur les côtés,
tous les sommets seraient sur $d_0$.

\textbf{Unimodularité.} « Pour des raisons techniques », on exige de plus que
les rotations soient réalisées par des automorphismes \emph{unimodulaires}
($\det u_V = 1$). Cela revient à demander que $\sigma_a$ (ou $\sigma_s$, qui
lui est conjuguée) soit une \emph{réflexion}, c'est-à-dire fixe une droite
point par point ; on comprendra la restriction en passant au plan projectif.
Les pages 82 à 87 justifient cette annonce : la condition exclut exactement
une famille, les réalisations « impropres », qui ne donnent de polygone fini
qu'en caractéristique $p \neq 2$, et ce sont alors des $2p$-gones réguliers non
unimodulaires\footnote{La page dit que la restriction « est automatiquement
satisfaite dans le cas des polygones réguliers », puis un mot illisible, puis « $2p$ côtés
en car.\ $p$ ($p$ nb premier impair) --- p.\ ex.\ l'hexagone en car.\ 3 ». Ce
qui est vrai, par les pages 84 et 85, est l'inverse pour cette famille : le
$2p$-gone impropre de caractéristique $p$ est régulier et \emph{n'est pas}
unimodulaire ($\det u_V = -1$), et il est le seul polygone régulier fini qui
ne le soit pas. Le mot illisible portait peut-être l'exception ; on ne le
restitue pas.}.

\textbf{Polygones propres, représentations géométriques.} Dans la définition
générale d'un polygone, les côtés géométriques sont des droites, et
l'application des côtés combinatoires vers les côtés géométriques n'est pas
forcément injective. Le polygone est \emph{propre} si un sommet $s$ est sur la
droite $\mathrm{dr}(s', s'')$ d'un côté combinatoire $\{s', s''\}$ seulement
quand $s \in \{s', s''\}$ ; les côtés géométriques sont alors en bijection avec
les côtés combinatoires. On verra que les polygones réguliers unimodulaires
sont propres (pages 101 à 105). La bonne notion, conclut la page 5, est celle
de \emph{représentation géométrique} d'un polygone combinatoire : deux
applications, l'une sur les sommets, l'autre vers les droites, qui respectent
l'incidence ; quand deux sommets adjacents ont des images distinctes, la
première détermine la seconde, et elle peut être choisie arbitrairement. Une
\emph{épinglage} d'un polygone géométrique par un polygone combinatoire $\Pi$
est un couple $(\Pi, \varphi)$, $\varphi : S \to E$\footnote{La page 5 est en
grande partie barrée et encadrée de renvois illisibles ; la phrase sur la
représentation géométrique se poursuit au-delà de la page. C'est la notion que
le cours reprend sous le nom de réalisation géométrique (page 50).}.

\pagerange{6}{9}
\subsection*{Les polygones réguliers réels (pages 6 à 9)}

Soit $E = \mathbf{C}$, vu comme plan réel, $\zeta$ une racine primitive
$n$-ième de l'unité ($n \geq 3$), $s_i = \zeta^i$ ($i \in
\mathbf{Z}/n\mathbf{Z}$) avec la structure polygonale standard, et $\sigma z =
\bar z$, $u z = \zeta z$ ; alors $\sigma s_i = s_{-i}$, $u s_i = s_{i+1}$, et
l'on obtient un $n$-gone régulier $P_\zeta$ --- convexe pour $\zeta =
e^{\pm 2i\pi/n}$, étoilé sinon. On vérifie, et la suite du dossier le
démontrera dans un cadre plus général, que :
\begin{enumerate}
  \item[a)] $P_\zeta \simeq P_{\zeta'}$ si et seulement si $\zeta' =
    \zeta^{\pm 1}$, ou encore si et seulement si $\alpha(\zeta) =
    \alpha(\zeta')$, où $\alpha(\zeta) = \zeta + \zeta^{-1} = 2\cos\arg\zeta$ ;
  \item[b)] tout polygone régulier d'un plan affine réel est isomorphe à un
    $P_\zeta$.
\end{enumerate}
Les classes d'isomorphie de polygones réguliers réels sont donc en bijection
avec l'ensemble des classes de racines de l'unité modulo l'inversion, privé
des classes de $1$ et de $-1$\footnote{La page écrit
$(\mathcal{T}/\{\pm 1\}) \smallsetminus \{\bar 1\}$, où $\mathcal{T} \simeq
\mathbf{Q}/\mathbf{Z}$ est le groupe des racines de l'unité, $\{\pm 1\}$
agissant par $\zeta \mapsto \zeta^{\pm 1}$. Il faut ôter aussi la classe de
$-1$ : $\zeta = -1$ donnerait un « polygone » à deux sommets. Ici
« isomorphisme » s'entend d'une bijection affine qui respecte la structure
combinatoire.}.

Pour quatre sommets consécutifs,
\[
  s_{i+3} - s_i = (1 + \alpha)\,(s_{i+2} - s_{i+1}), \qquad
  \frac{\zeta^{i+3} - \zeta^i}{\zeta^{i+2} - \zeta^{i+1}}
  = \frac{\zeta^2 + \zeta + 1}{\zeta} = 1 + \zeta + \zeta^{-1} ,
\]
et le facteur réel $1 + \alpha$ ne dépend que de $\zeta$. Cette relation, qui
ne fait intervenir que des rapports de vecteurs parallèles, est affine ; elle
reviendra comme définition (page 82).

Si l'on étend les scalaires à $\mathbf{C}$, $E \otimes_{\mathbf{R}} \mathbf{C}
\simeq \mathbf{C}^2$ par $z \mapsto (z, \bar z)$, la multiplication par
$\zeta$ se diagonalise suivant les deux droites isotropes,
$u = \mathrm{diag}(\zeta, \zeta^{-1})$, le sommet $s_i$ devient $(\zeta^i,
\zeta^{-i})$ et la symétrie $z \mapsto \bar z$ l'échange des coordonnées.
C'est, annonce la page 8, la structure typique des $n$-gones réguliers sur un
corps algébriquement clos de caractéristique première à $n$ (elle est établie
page 63), avec la même classification que sur $\mathbf{R}$ ; et il est
« remarquable », note-t-il, que $\mathbf{R}$, qui n'est pas algébriquement
clos, donne la même classification. L'explication viendra aux pages 81
à 90 : un polygone régulier est déterminé par $\alpha$ seul, sans qu'il soit
besoin que $\zeta$ soit dans le corps, et sur $\mathbf{R}$ toutes les valeurs
$2\cos(2\pi k/n)$ sont réelles.

La page 9 est une feuille de calculs isolée, barrée, sur les endomorphismes
d'un plan : pour $u$, $v$ d'ordre $2$,
\[
  \det(u + v) = \det u + \det v + \operatorname{Tr} u \operatorname{Tr} v -
  \operatorname{Tr}(uv), \qquad
  \operatorname{Tr}(u \check v) = \operatorname{Tr} u \operatorname{Tr} v -
  \operatorname{Tr}(uv),
\]
où $\check v = (\operatorname{Tr} v)\,\mathrm{id} - v$ est l'adjoint
classique ; puis des essais sur les endomorphismes de trace nulle laissant
stable une droite, nilpotents ou de carré nul\footnote{Les deux identités sont
justes en dimension $2$. Le reste de la feuille (des fragments où paraissent
$\mathcal{O}_\xi \subset \sqrt{\mathcal{J}} \subset \mathcal{L}$ et une
flèche $X \to \mathbf{R}(\sqrt{\mathcal{J}}/\mathcal{O}_\xi)$) est trop
lacunaire pour être lu ; rien ne le rattache aux pages voisines.}.

\pagerange{10}{11}
\subsection*{Repères et épinglages (pages 10 et 11)}

Un \emph{repère} d'un polygone combinatoire $\Pi = (S, \omega)$ est un couple
$(s_0, d_0)$ formé d'un sommet et d'un côté incidents --- un \emph{drapeau},
dirait-on aujourd'hui --- ; il revient au même de se donner deux sommets
adjacents $(s_0, s_1)$, ou un sommet $s_0$ et une orientation $u \in \omega$
($s_1 = u s_0$).

\textbf{Proposition 1.} Le groupe $\Gamma = \mathrm{Aut}(S, \omega)$ opère
simplement transitivement sur l'ensemble $\mathrm{Rep}(\Pi)$ des repères :
c'est un torseur sous $\Gamma$.

\textbf{Corollaire 1.} Si $\Pi$ et $\Pi'$ sont isomorphes et $r_0$ est un
repère de $\Pi$, l'application $\mathrm{Isom}(\Pi, \Pi') \to
\mathrm{Rep}(\Pi')$, $\theta \mapsto \theta(r_0)$, est bijective.

Un \emph{épinglage} d'un polygone combinatoire est le choix d'un repère ;
c'est la même chose qu'un isomorphisme avec le $n$-gone combinatoire type
$\Pi_n$ muni de son repère standard. Un épinglage d'un polygone géométrique
est un épinglage du polygone combinatoire sous-jacent.

\textbf{Corollaire 2.} Si $P$ et $P'$ sont deux polygones réguliers
isomorphes et $r$ un repère de $P$, l'application $\mathrm{Isom}(P, P') \to
\mathrm{Rep}(P')$, $\theta \mapsto \theta(r)$, est bijective.

\textbf{Scholie.} Dans une classe d'isomorphie donnée, les $n$-gones réguliers
épinglés sont déterminés à isomorphisme \emph{unique} près, et les classes
d'isomorphie de $n$-gones épinglés et non épinglés se correspondent
bijectivement. La page 11 pose alors la « question » qui organise la suite :
écrire « le » plan $E$ où vit $P$, et les sommets $s_i$, en fonction du type
d'isomorphie, c'est-à-dire de $n$ et d'un « module » $\alpha \in k$. La
réponse sera la formule $s_n = (F_n F_{n+1}(\alpha), F_{n-1} F_n(\alpha))$
des pages 111 et 112.

\pagerange{12}{21}
\subsection*{Polygones réguliers et représentations du groupe diédral (pages 12 à 21)}

Un polygone régulier $P = (S, \omega)$ détermine un sous-groupe $\Gamma
\subset \mathrm{Aff}(E)$, isomorphe à $\mathbf{D}_n$ (sans isomorphisme
privilégié), dont $S$ est une orbite. Le groupe ne suffit pas à retrouver
$\omega$ : $\omega$ et $\omega^{(r)} = \{u^r, u^{-r}\}$, $(r, n) = 1$, ont le
même groupe d'automorphismes. Il faut donc se donner en plus la partie
$\omega \subset \Gamma$, avec les conditions
\begin{enumerate}
  \item[(6)] le sous-groupe $\Gamma^{+}$ engendré par $\omega$ est d'indice
    $2$, $\Gamma/\Gamma^{+}$ opère sur $\Gamma^{+}$ par $\lambda \mapsto
    \lambda^{-1}$, et l'extension est scindée par une involution $\sigma \in
    \Gamma \smallsetminus \Gamma^{+}$ ;
  \item[(7)] $\Gamma^{+}$ est fini ;
  \item[(8)] $S$ est un torseur sous $\Gamma^{+}$ ;
  \item[(9)] $S$ n'est pas contenu dans une droite.
\end{enumerate}
Ces conditions suffisent : la structure de torseur sous le groupe cyclique
$\Gamma^{+}$ fait de $S$ un polygone combinatoire, que $\sigma$ retourne, et
$\Gamma$ s'injecte dans son groupe d'automorphismes avec $\Gamma^{+}
\xrightarrow{\sim} \mathrm{Aut}^{+}$\footnote{Le texte qui justifie la
suffisance (page 14) est en partie illisible ; on en donne l'argument. La page
numérote les conditions dans l'ordre (9), (8), (7), et des flèches dans la
marge rétablissent (7), (8), (9). En (7$'$), « $\operatorname{card}
\Gamma^{+} \simeq \mathbf{Z}/n\mathbf{Z}$ » est à lire $\Gamma^{+} \simeq
\mathbf{Z}/n\mathbf{Z}$.}. On peut remplacer la donnée de $S$ par celle d'une
orbite sous $\Gamma^{+}$ et de la condition
\[
  (10) \qquad \exists\, \sigma_{s_0} \in \Gamma,\ \sigma_{s_0} \neq 1,\
  \sigma_{s_0} s_0 = s_0 ,
\]
ce $\sigma_{s_0}$ étant a posteriori unique.

Si l'on épingle entièrement (un sommet $s_0$ et une orientation $u$), les
données deviennent, pour $n \geq 3$ fixé, un homomorphisme $\varphi :
\mathbf{D}_n \to \mathrm{Aff}(E)$ et un point $s_0 \in E$, soumis à
\[
  (13) \qquad \sigma_1 s_0 = s_0, \qquad s_0,\ s_1 = u s_0,\ s_2 = u^2 s_0
  \ \text{non alignés},
\]
\[
  (12) \qquad \varphi \text{ fidèle}, \qquad u \text{ unimodulaire}.
\]

\textbf{Proposition 2.} Pour un plan $E$ et un entier $n \geq 3$ fixés, se
donner dans $E$ un $n$-gone régulier (unimodulaire) épinglé revient à se donner
un couple $(\varphi, s_0)$ comme ci-dessus.

La donnée essentielle est $\varphi$. On verra que, pour les $\varphi$ qui
proviennent de polygones réguliers, $\sigma_1$ est une réflexion : son lieu
fixe $D_{\sigma_1}$ est une droite, et le choix de $s_0$ est celui d'un point
de cette droite. Comme $u s_0 = \sigma_0 \sigma_1 s_0 = \sigma_0 s_0$, la
condition $u s_0 \neq s_0$ demande $s_0 \notin D_{\sigma_0}$ ; on trouvera que
$\sigma_0$ est aussi une réflexion et que $D_{\sigma_0} \neq D_{\sigma_1}$,
ce qui exclut au plus un point, le \emph{centre} $c$, point fixe de tout
$\varphi(\mathbf{D}_n)$. Soit $D_{\sigma_1}^{*} = D_{\sigma_1} \smallsetminus
(D_{\sigma_0} \cap D_{\sigma_1})$, sur lequel opère le commutant $\mathfrak{G}$
de $\varphi(\mathbf{D}_n)$ dans $\mathrm{Aff}(E)$.

\textbf{Proposition 3.} Supposons que $\sigma_0$ et $\sigma_1$ soient des
réflexions de droites fixes distinctes --- ce qui, note-t-il, n'est autre que
la condition unimodulaire.
\begin{enumerate}
  \item[1°)] Si $D_{\sigma_0}$ et $D_{\sigma_1}$ se coupent en $c$, alors $c$
    est l'unique point fixe de $\Gamma$ ; prenant $c$ pour origine, $\Gamma$
    opère linéairement, et $\mathfrak{G}$ est le groupe des homothéties de
    centre $c$, $\mathfrak{G} \simeq k^{*}$.
  \item[2°)] Si $D_{\sigma_0} \parallel D_{\sigma_1}$, soit $T \simeq k$ le
    groupe des translations de leur direction commune ; alors $\mathfrak{G} =
    T$.
\end{enumerate}
Dans les deux cas, $D_{\sigma_1}^{*}$ est un torseur sous
$\mathfrak{G}$\footnote{La démonstration de la page 19 : le groupe évident
$\mathfrak{G}_0$ ($k^{*}$, resp.\ $T$) commute à $\sigma_0$ et $\sigma_1$,
donc à $\Gamma$, et opère simplement transitivement sur $D_{\sigma_1}^{*}$ ;
$\mathfrak{G}$ opère librement sur l'ensemble des $s_0$ admissibles (un
élément du commutant qui fixe $s_0$ fixe tous les $s_i$, qui engendrent $E$) ;
il en résulte « formellement » que $\mathfrak{G} = \mathfrak{G}_0$. Il faut
que $\varphi$ soit admissible, c'est-à-dire qu'il existe un $s_0$ satisfaisant
(13) ; alors, $\mathfrak{G}_0$ étant transitif et préservant (13), tous les
points de $D_{\sigma_1}^{*}$ le sont. L'indice de « $\mathfrak{G}_0 =
\mathfrak{G}$ » est lu avec doute. Par les pages 87 et 88, le cas 2° est
celui de $\alpha = 2$, qui ne donne de polygone fini qu'en caractéristique
$p$, avec $n = p$ ; la marge de la page 20, en partie illisible, semble le
dire (« $n \neq \operatorname{car} k$ \ldots{} $n \neq 2$ si $p = 2$
»).}.

\textbf{Corollaire 1 (page 20).} Deux $n$-gones réguliers épinglés dont les
représentations de $\mathbf{D}_n$ sont isomorphes sont isomorphes : classer
les représentations affines admissibles de $\mathbf{D}_n$, c'est classer les
$n$-gones réguliers.

\textbf{Corollaire 2 (pages 20 et 21).} Si $\varphi$, $\varphi'$ sont deux
représentations admissibles et $P$ un $n$-gone régulier épinglé subordonné à
$\varphi$, l'application $\theta \mapsto \theta(P)$ est une bijection de
$\mathrm{Isom}((E, \varphi), (E', \varphi'))$ sur l'ensemble des $n$-gones
réguliers épinglés subordonnés à $\varphi'$\footnote{La page 21 ajoute
« ((19) disait que c'était une surjection) », renvoi qu'on ne retrouve pas ;
elle écrit $D_n^{*}$ pour $D_{\sigma_1}^{*}$, et $\square$ pour
$\mathbf{D}_n$ au corollaire 1.}.

En langage suggestif : choisir un sommet $s_0 \in D_{\sigma_1}^{*}$, c'est
\emph{épingler} la représentation $\varphi$ --- on tue ses automorphismes
($\mathfrak{G}$ opère librement) sans introduire de nouvelles classes
d'isomorphie ($\mathfrak{G}$ opère transitivement).

\pagerange{22}{26}
\subsection*{Où l'on laisse tomber la fidélité (pages 22 à 26)}

Considérons des données $(\varphi, s_0)$ qui satisfont (13), mais pas
nécessairement (12) : ni la fidélité, ni l'unimodularité.

\textbf{Proposition 4.}
\begin{enumerate}
  \item[(a)] Pour $v \in \mathbf{D}_n^{+}$, $\varphi(v) = \mathrm{id}_E$ si et
    seulement si $v s_0 = s_0$.
  \item[(b)] $\varphi$ est fidèle si et seulement si sa restriction à
    $\mathbf{D}_n^{+}$ l'est, si et seulement si $u^d s_0 \neq s_0$ pour tout
    diviseur \emph{strict} $d$ de $n$.
  \item[(c)] Le noyau de $\varphi$ est le sous-groupe engendré par
    $u^{n'}$, où $n' = \operatorname{card} S$ est un diviseur de $n$, et
    $n' \geq 3$ ; $\varphi$ se factorise par une représentation fidèle
    $\varphi'$ de $\mathbf{D}_{n'}$\footnote{Variances de la page. En (b) :
    « $\forall d \mid n$, on ait $u^d(s_0) \neq s_0$ », où il faut exclure
    $d = n$ ; un signe isolé sous $\mathbf{N}^{*}$, non lu, était peut-être
    cette restriction. En (c) : « le noyau est engendré par un élément $u^d$,
    où $d$ est un diviseur de $n$ tel que $n' = n/d \geq 3$ », puis, dans la
    démonstration, $n = n'd$ et $\operatorname{card} S = n'$. Les deux ne
    s'accordent pas : si le noyau est $\langle u^d \rangle$, l'image de
    $\mathbf{D}_n^{+}$ est d'ordre $d$ et $\operatorname{card} S = d$. On
    énonce ce qui est vrai, avec $n'$ le cardinal de $S$. La démonstration
    de (b) n'est pas écrite ; elle résulte de (a) et du corollaire 1.}.
\end{enumerate}
\emph{Démonstration.} (a) Les $s_\nu = u^\nu s_0$ engendrent affinement $E$, et
$v s_\nu = u^\nu v s_0$ puisque $\mathbf{D}_n^{+}$ est commutatif. (c) Le noyau
de $\varphi|_{\mathbf{D}_n^{+}}$ est un sous-groupe de $\mathbf{Z}/n$, engendré
par $u^{n'}$ avec $n' \mid n$ ; alors $\operatorname{card} S = n'$, et $n'
\geq 3$ parce que $S$ engendre $E$. $\varphi$ se factorise par $\varphi' :
\mathbf{D}_{n'} \to \mathrm{Aff}(E)$, fidèle sur les rotations, donc fidèle
par (b), et le noyau de $\varphi$ est celui de sa restriction aux rotations.

\textbf{Corollaire 1.} Si $v \in \mathbf{D}_n \smallsetminus
\mathbf{D}_n^{+}$, $\varphi(v) \neq \mathrm{id}_E$ : $\varphi(v)$ échange les
deux ordres circulaires de $S$ ($\operatorname{card} S \geq 3$), donc n'est
pas l'identité sur $S$.

\textbf{Corollaire 2.} Dans la situation de (c), $(\varphi', s_0)$ définit
dans $E$ un $n'$-gone épinglé.

\textbf{Scholie.} La donnée $(\varphi, s_0)$ est une réalisation géométrique
du $n$-gone combinatoire $\Pi_n$ qui est un revêtement à $n/n'$ feuillets de
la réalisation $(\varphi', s_0)$ de $\Pi_{n'}$.

Quand on fixe $n$, on trouve d'abord des représentations $\varphi$ sans savoir
lesquelles sont fidèles, et c'est une question d'un autre type ; surtout sur
un anneau plus général qu'un corps --- $\mathbf{Z}[T]/(T^n - 1)$, les anneaux
cyclotomiques ---, où certaines fibres donnent des représentations fidèles et
d'autres non, sans que celles-ci cessent de définir des polygones réguliers,
à moins de côtés. D'où la règle de conduite :
\begin{quote}
« Ce serait emprisonnant --- et finalement d'ailleurs artificiel --- de
traîner tout au long dans les données la condition de fidélité, au lieu de se
borner à l'examiner seulement aux moments où cette considération est
utile\footnote{Citation de la page 26 ; « finalement d'ailleurs » est un
ajout lu avec doute. La phrase sur les fibres (page 25) est lacunaire ; on en
donne le sens que les pages 114 à 119 confirment.}. »
\end{quote}

\pagerange{26}{33}
\subsection*{Les $\infty$-polygones (pages 26 à 33)}

On peut aller plus loin. La représentation $\varphi$ revient à deux
involutions $\sigma_0, \sigma_1 \in \mathrm{Aff}(E)$ avec $(\sigma_0
\sigma_1)^n = \mathrm{id}$, et la fidélité s'écrit
\[
  (15) \qquad u^d s_0 \neq s_0 \quad \text{pour tout diviseur strict } d
  \text{ de } n .
\]
Comme $s_0, s_1, s_2$ sont distincts, les diviseurs $d \leq 2$ sont sans
objet : il suffit de tester les diviseurs stricts $d \geq 3$. La condition
est vide exactement quand $n$ est un nombre premier impair ou $n = 4$ ; le
premier cas où elle ne l'est pas est $n = 6$, où l'hexagone régulier peut
dégénérer en un triangle parcouru deux fois\footnote{La page dit : « il
suffit de l'exiger pour $3 \leq d \leq [\frac n3]$ (donc si $3 \leq n \leq 8$,
cette condition est automatiquement satisfaite, le premier entier où elle ne
l'est pas est $n = 9$ \ldots{} le $9$-gone régulier pouvant dégénérer en trois
fois un $3$-gone régulier) ». Si $d$ y désigne le nombre de feuillets
$n/n'$, la borne $d \leq n/3$ est juste ($n' \geq 3$), mais la borne
inférieure doit être $d \geq 2$ : le triangle parcouru deux fois (sur un corps
de caractéristique $\neq 3$, $\alpha = -1$, $u^3 = 1$, donc $u^6 = 1$) est
une représentation non fidèle de $\mathbf{D}_6$ satisfaisant (13), de même
que le carré parcouru deux fois pour $n = 8$. Avec la borne corrigée, la
condition est vide pour $n$ premier impair et pour $n = 4$, et pour eux
seulement ; c'est ce que dira, pour $n = p$, la page 115.}.

Mais on peut aussi laisser tomber $(\sigma_0\sigma_1)^n = \mathrm{id}$
lui-même, c'est-à-dire regarder une représentation
\[
  (16) \qquad \varphi : \mathbf{D}_\infty \longrightarrow \mathrm{Aff}(E)
\]
du groupe diédral infini, en s'intéressant surtout à celles qui se factorisent
par un $\mathbf{D}_n$. De ce point de vue la fidélité ne se pose plus : les
représentations les plus intéressantes de $\mathbf{D}_\infty$, celles qui
donnent les polygones finis, sont précisément non fidèles.

Il y a lieu alors de considérer des polygones combinatoires infinis, et leurs
réalisations géométriques \emph{régulières} : une application $\varphi_S : S
\to E$ telle que
\[
  (17) \qquad \forall v \in \mathrm{Aut}(\Pi),\ \exists!\, v_E \in
  \mathrm{Aff}(E),\quad v_E \circ \varphi_S = \varphi_S \circ v .
\]
Cela revient à un homomorphisme $\varphi_\Gamma : \Gamma = \mathrm{Aut}(\Pi)
\to \mathrm{Aff}(E)$ et une application $\Gamma$-équivariante $\varphi_S : S
\to E$ dont l'image n'est pas contenue dans une droite (19). Si l'on fixe un
sommet $i_0$ et l'unique élément $\neq 1$ de $\Gamma$ qui le fixe\footnote{La
page l'appelle $\sigma_0$ ; c'est notre $\sigma_1$.}, la donnée de
$\varphi_S$ équivaut à celle de $s_0 = \varphi_S(i_0)$ fixé par cet élément,
et (19) à la non-colinéarité de $s_0, s_1, s_2$ (22). Se donner $(s_0, u)$,
c'est épingler $\Pi$ par $\Pi_n$, $3 \leq n \leq \infty$.

\textbf{Proposition 5.} Les réalisations géométriques régulières de
$\Pi_\infty$ dans $E$ correspondent aux triples $(\sigma_0, \sigma_1, s_0)$ de
deux automorphismes affines et d'un point tels que
\begin{enumerate}
  \item[1°)] $\sigma_0^2 = \sigma_1^2 = \mathrm{id}$ ;
  \item[2°)] $\sigma_1 s_0 = s_0$ ;
  \item[3°)] $s_0$, $s_1 = u s_0$, $s_2 = u^2 s_0$ ne sont pas alignés,
    $u = \sigma_0 \sigma_1$.
\end{enumerate}
La réalisation est \emph{unimodulaire} si de plus 4°) $u$ est unimodulaire.

\textbf{Scholie.} On étudie ainsi d'un coup tous les $n$-gones réguliers
épinglés --- ceux où $u$ est d'ordre fini $n$ --- et un objet nouveau, les
$\infty$-polygones réguliers épinglés.

\textbf{Remarque.} Les résultats précédents (propositions 1 à 4 et leurs
corollaires) s'étendent sans dommage aux $\infty$-polygones, la divisibilité
de la proposition 4 perdant son sens pour $n = \infty$. La marge ajoute qu'il
conviendrait de \emph{définir} d'emblée les polygones réguliers généraux,
finis ou non --- ce que fera le cours.

\textbf{NB (pages 32 et 33).} Le module $\alpha \in k$ n'est plus soumis à
aucune condition, cyclotomique ou autre. Plus intrinsèquement, pour un polygone
combinatoire $\Pi$ fini ou infini, la catégorie de ses réalisations
géométriques régulières est \emph{rigide} : une réalisation n'a pas
d'automorphisme autre que l'identité (un automorphisme compatible fixe $S$
point par point, et $S$ engendre $E$). On peut donc parler de \emph{la}
réalisation $E(\Pi)$ de module $\alpha$, sur laquelle l'action de
$\mathrm{Aut}(\Pi)$ s'obtient automatiquement --- un des intérêts techniques
d'éliminer l'épinglage dès le départ\footnote{Les deux pages sont rapides et
lacunaires ; le mot « rigide » et la parenthèse « sans automorphismes autres
que l'identité » sont sur la page, la justification est de nous.}.

\pagerange{33}{39}
\subsection*{L'étoilage (pages 33 à 39)}

Soit $(S, \omega)$ un polygone combinatoire, $\omega = \{u, u^{-1}\}$, et $r
\in \mathbf{Z}$ ; posons $\omega^{(r)} = \{u^r, u^{-r}\}$. Passer de $\omega$
à $\omega^{(r)}$, c'est relier chaque sommet à son $r$-ième voisin : le
pentagone devient le pentagramme pour $r = 2$.

\textbf{Proposition 6.}
\begin{enumerate}
  \item[a)] $(S, \omega^{(r)})$ est un polygone combinatoire si et seulement si
    $n = \operatorname{card} S$ est fini et $(r, n) = 1$, ou $r = \pm 1$.
  \item[b)] Si $n$ est fini, $\omega^{(r)} = \omega^{(r')}$ si et seulement si
    $r' \equiv \pm r \pmod n$.
  \item[c)] $(\omega^{(r')})^{(r)} = \omega^{(rr')}$. Le groupe $Q_n =
    (\mathbf{Z}/n\mathbf{Z})^{*}/\{\pm 1\}$ opère donc, et librement, sur
    l'ensemble $\mathrm{Pol}(S)$ des structures polygonales sur $S$.
\end{enumerate}
\emph{Démonstration.} Une structure polygonale fait de $S$ un torseur sous
$\mathbf{Z}_n$ ($\mathbf{Z}/n\mathbf{Z}$ si $n$ est fini, $\mathbf{Z}$ sinon) ;
l'action déduite de $u^r$ se déduit de la précédente par la multiplication par
$r$ de $\mathbf{Z}_n$, et fait de $S$ un torseur si et seulement si celle-ci
est bijective. b) $\omega^{(r)} = \omega^{(r')}$ signifie $u^{r'} = u^{\pm
r}$\footnote{La page conclut a) par « $n$ fini et $(r, n) = 1$ » : pour $n$
infini, la multiplication par $r = \pm 1$ est bijective, et $\omega^{(\pm 1)}
= \omega$. L'alinéa marqué (c) de la démonstration prouve b).}.

\textbf{Remarque.} Une structure de polygone \emph{orienté} est la donnée d'une
permutation circulaire $u$, c'est-à-dire d'une structure de
$\mathbf{Z}_n$-torseur ; $(\mathbf{Z}/n\mathbf{Z})^{*}$ opère librement sur
l'ensemble $\mathrm{Polor}(S)$ de ces structures, par $u \mapsto u^r$.

Si $P = (E, S, \omega)$ est un $n$-gone régulier et $(r, n) = 1$,
$P^{(r)} = (E, S, \omega^{(r)})$ est encore un $n$-gone régulier, de même
ensemble de sommets : les normalisateurs de $\omega$ et de $\omega^{(r)}$ dans
$\mathfrak{S}_S$ sont les mêmes,
\[
  (27) \qquad \mathrm{Aut}(\Pi) = \mathrm{Aut}(\Pi^{(r)}) \subset
  \mathfrak{S}_S ,
\]
et l'unimodularité passe de $u$ à $u^r$. $Q_n$ opère donc librement sur
l'ensemble des $n$-gones réguliers de $E$, et $(\mathbf{Z}/n)^{*}$ sur celui
des $n$-gones réguliers orientés, ou orientés et munis d'un
sommet\footnote{En (26), la page écrit $\Psi_r(E)$ pour ce qui est ici
$P^{(r)}$.}.

\textbf{Question} (pages 36 et 37). Un $n$-gone régulier $P$ peut-il être
isomorphe à $P^{(r)}$ pour $r \not\equiv \pm 1$ ? En termes de l'invariant,
$P^{(r)}$ a pour invariant $S_r(\alpha) = \zeta^r + \zeta^{-r}$, et
l'isomorphie équivaut à $S_r(\alpha) = \alpha$. Si la caractéristique $p$ ne
divise pas $n$, $\zeta$ est d'ordre exactement $n$ et $\zeta^r \neq
\zeta^{\pm 1}$ : jamais. Si $p \mid n$, les seuls $n$-gones réguliers sont,
par les pages 94, 95 et 118, les $p$-gones d'invariant $2$ et les
$2p$-gones d'invariant $-2$ (dont le carré en caractéristique $2$) ; il n'y a
alors qu'une classe, et $P^{(r)} \simeq P$ pour tout $r$\footnote{Le passage
des pages 36 et 37 est très rapide, plusieurs fois récrit, et seules les
affirmations mathématiques y sont sûres : « si $(n, p) = 1$ on aura
$\Psi_r(\alpha) \neq \alpha$ si $r \neq 1$ dans $\Psi_n$, tandis que si $p
\mid n$ \ldots{} (ou bien $n = p$, $p$ impair, ou bien $n = 2p$) alors
\ldots{} les $n$-pol.\ réguliers sont non isom. », la lecture « non »
étant douteuse. Ce qu'on énonce est ce que les pages 94 à 119 établissent. La
suite, sur l'ordre de $r$ dans $(\mathbf{Z}/p)^{*}$ et « l'hexagone », n'a pas
pu être suivie. $S_r$ est le polynôme de la page 122.}. On retiendra, dit-il,
que l'étoilage « a une nette tendance à changer la classe d'isomorphie ».

Pour les représentations de $\mathbf{D}_\infty$, l'étoilage se définit pour
\emph{tout} $r \in \mathbf{Z}$ : on remplace $(\sigma_1, u)$ par $(\sigma_1,
u^r)$, ce qui conserve les relations $\sigma_1^2 = 1$, $\sigma_1 u^r
\sigma_1 = u^{-r}$. Mais la non-colinéarité de $s_0, u^r s_0, u^{2r} s_0$ peut
tomber en défaut : si la représentation se factorise par $\mathbf{D}_n$ et $r
= n$, tous les points sont confondus. La page pose alors la question de la
signification géométrique de la représentation de $\mathbf{D}_\infty$
attachée à l'invariant étoilé $S_r(\alpha)$ dans ces cas « exceptionnels » où
les $s_i$ étoilés sont alignés\footnote{La question est sur la page ; elle
reste sans réponse. On remarquera (c'est de nous) que dans le cas $\zeta^r =
\pm 1$, l'invariant étoilé est $S_r(\alpha) = \pm 2$, celui du polygone
parabolique ou de la bande des pages 94 à 100 : la représentation propre
d'invariant $S_r(\alpha)$ n'est donc pas la représentation étoilée dégénérée,
et c'est sans doute ce que la question vise.}.

\pagerange{39}{48}
\subsection*{Structures quasi-polygonales et systèmes réguliers de sommets (pages 39 à 48)}

Une \emph{structure quasi-polygonale} sur un ensemble fini $S$ de cardinal $n$
est une orbite $\pi$ de $Q_n$ dans $\mathrm{Pol}(S)$ --- une structure
polygonale « modulo étoilage ». Le sous-groupe $Z \subset \mathfrak{S}_S$
engendré par un $\omega \in \pi$ ne dépend que de $\pi$ ; il est cyclique
d'ordre $n$ et $S$ est un $Z$-torseur. Se donner une structure
quasi-polygonale, c'est donc se donner une structure de torseur sous un groupe
cyclique d'ordre $n$ ; choisir $\omega \in \pi$, c'est choisir un générateur
de $Z$ au signe près ; choisir une orientation, c'est choisir un générateur.

En termes d'aujourd'hui, $Z$ est un $\mathbf{Z}/n\mathbf{Z}$-module libre de
rang $1$ et $S$ un torseur sous lui : une \emph{droite affine sur
$\mathbf{Z}/n\mathbf{Z}$}. La catégorie des quasi-polygones combinatoires
d'ordre $n$ est équivalente à celle des fibrés affines de rang $1$ sur
$\mathrm{Spec}(\mathbf{Z}/n\mathbf{Z})$\footnote{La page 40 écrit « fibrés
» (un mot illisible) « affines sur $\mathrm{Spec}(\mathbf{Z}/n\mathbf{Z})$ », et la
page 41 « l'ensemble de ses sections, qui est un torseur sous le groupe des
sections du fibré vectoriel correspondant, lequel est un
$\mathbf{Z}/n\mathbf{Z}$-module libre de rang $1$ ».}. Son groupe
d'automorphismes $\mathrm{Aff}(S)$ est le groupe affine de cette droite :
\[
  0 \longrightarrow Z \longrightarrow \mathrm{Aff}(S) \longrightarrow
  (\mathbf{Z}/n\mathbf{Z})^{*} \simeq \mathrm{Aut}(Z) \longrightarrow 0 ,
\]
extension scindée par le choix d'un point $s \in S$, qui identifie
$\mathrm{Aff}(S)$ au produit semi-direct $(\mathbf{Z}/n)^{*} \ltimes
\mathbf{Z}/n$. L'image inverse $\mathrm{Aff}^{\pm}(S)$ de $\{\pm 1\}$ est le
groupe diédral $\mathrm{Aut}(S, \omega)$ de toute structure $\omega \in \pi$ ;
il est d'indice $\varphi(n)/2$ dans $\mathrm{Aff}(S)$, et lui est égal
exactement quand $\varphi(n) = 2$, c'est-à-dire $n = 3, 4, 6$\footnote{La page
dit $\mathrm{Aff}^{\pm}(S)$ « d'indice deux » (lecture incertaine) ; l'indice
est $[(\mathbf{Z}/n)^{*} : \{\pm 1\}] = \varphi(n)/2$, qui vaut $2$ pour
$n = 5, 8, 10, 12$ seulement. Le crochet de la page 42 sur $n = 3, 4, 6$
(« Vérifier ») est juste.}.

\textbf{Systèmes réguliers de sommets.} Un \emph{système régulier de $n$
sommets} dans $E$ (la page abrège SR$n$S) est un couple $(S, Z)$, $S \subset
E$ finie de cardinal $n$, $Z$ une structure quasi-polygonale sur $S$, telle que
pour une (donc toute) structure polygonale $\omega$ compatible, $(E, S,
\omega)$ soit un polygone régulier --- « donc toute » par (27). C'est la même
chose qu'une orbite de $Q_n$ opérant sur les $n$-gones réguliers de $E$.

\textbf{Question} (pages 43 et 44) : $Z$ est-il déterminé par la seule partie
$S \subset E$ ? Le stabilisateur de la partie $S$ dans $\mathrm{Aff}(E)$ est
$\mathrm{Aut}(E, S, \omega) \simeq \mathbf{D}_n$, sauf pour $\alpha = \pm 2$,
où c'est $\mathrm{Aff}(S, Z)$ tout entier\footnote{Énoncé de la page,
marqué « vérifier ! ». On le vérifie (c'est de nous). Hors des cas $\alpha =
\pm 2$, la caractéristique ne divise pas $n$, les sommets centrés vérifient
$s_{i+1} + s_{i-1} = \alpha\, s_i$, et une application affine qui envoie
$s_i$ sur $s_{ri}$ imposerait $S_r(\alpha) = \alpha$, donc $r \equiv \pm 1$.
Pour $\alpha = 2$ (le $p$-gone de la page 100, sommets $(i, i^2)$ dans les
coordonnées $(v, w)$), $i \mapsto ri$ est induit par l'application linéaire
$(v, w) \mapsto (rv, r^2 w)$. Pour $\alpha = -2$ (le $2p$-gone de la
page 95, sommets pairs $s_0 + m\,a$ et impairs $s'_0 - m\,a$ sur deux droites
parallèles), $i \mapsto ri$ ($r$ impair) est induit par l'application affine
qui multiplie par $r$ le long de la direction $a$ et envoie $s'_0$ sur $s'_0 -
\frac{r-1}{2}\,a$.}. La question devient combinatoire : $Z$ est-il déterminé
par $\mathrm{Aff}^{\pm}(S) \subset \mathfrak{S}_S$, ou par $\mathrm{Aff}(S)
\subset \mathfrak{S}_S$ ? Dans le premier cas oui : $Z$ est le plus grand
sous-groupe cyclique distingué de $G = \mathrm{Aff}^{\pm}(S) \simeq
\mathbf{D}_n$ distinct de $G$, car le sous-groupe distingué engendré par une
réflexion $\sigma \in G \smallsetminus Z$ n'est pas cyclique\footnote{La page
dit qu'il est isomorphe à $\{\pm 1\} \cdot \mathbf{Z}/n'$ avec $n' = n/2$ ;
c'est vrai pour $n$ pair (pour $n = 4$, c'est le groupe de Klein) ; pour $n$
impair il est $G$ tout entier. Dans les deux cas il n'est pas cyclique.}. Le
cas de $\mathrm{Aff}(S)$ est laissé (« faire dém. »)\footnote{Il se traite
aussi (de nous) : ce cas n'arrive que pour $n = p$ ou $2p$ en caractéristique
$p$, et dans $\mathrm{Aff}(\mathbf{Z}/p) = (\mathbf{Z}/p)^{*} \ltimes
\mathbf{Z}/p$, comme dans $\mathrm{Aff}(\mathbf{Z}/2p) \simeq \mathbf{Z}/2
\times \mathrm{Aff}(\mathbf{Z}/p)$, les sous-groupes distingués cycliques sont
contenus dans les translations.}. « Donc sauf erreur », la structure
quasi-polygonale d'un système régulier de sommets est unique.

\textbf{Principe de classification géométrique.} Pour $n$ fixé, un système
régulier de $n$ sommets détermine (a) un sous-groupe cyclique $Z \subset
\mathrm{Aff}(E)$ d'ordre $n$, (b) une orbite $S$ de $Z$ dans $E$, soumis à
\begin{enumerate}
  \item[1)] $S$ n'est pas contenu dans une droite (donc $Z$ opère fidèlement
    sur $S$, qui est un $Z$-torseur) ;
  \item[2)] il existe $\sigma \neq 1$ dans le normalisateur de $Z$, qui opère
    sur $Z$ par $\lambda \mapsto \lambda^{-1}$ et fixe un point de $S$.
\end{enumerate}

\textbf{Proposition 7.} Se donner un système régulier de $n$ sommets dans $E$
équivaut à se donner $(Z, S)$ comme en (a), (b), satisfaisant 1) et 2)
(et, si l'on tient à l'unimodularité du chapitre, $Z$ engendré par un élément
de déterminant $1$)\footnote{La page a barré, dans 2), « opère sur $Z$ par
$\lambda \mapsto \lambda^{-1}$ » et ne garde que « $\exists\, \sigma \in
\mathrm{Norm}(Z)$ tel que $\sigma$ laisse invariant un pt de $S$ » ; sous
cette forme la condition est remplie par $\sigma = 1$ et ne dit rien. On
rétablit la clause barrée et $\sigma \neq 1$, qui sont ce qu'il faut pour que
$\{1, \sigma\} \cdot Z$ (produit semi-direct, page 45) opère sur $S$ en
groupe diédral. La proposition ajoute « éventuellement la condition de
minimalité : 3) $Z$ formé d'ordre minimal », « formé » lu avec doute,, qu'on n'a pas su
interpréter.}.

\textbf{Remarque.} Selon $p = \operatorname{car} k$ et $n$, le groupe des
automorphismes d'un système régulier de $n$ sommets est $\mathrm{Aff}^{\pm}(S)
\simeq \mathbf{D}_n$ ou $\mathrm{Aff}(S)$.

\textbf{Un programme de descente} (pages 46 à 48). On cherche une notion
d'épinglage pour ces systèmes, qui ramène leur étude à un calcul ; la théorie
de la descente devrait alors fournir un objet épinglé canonique sur le corps
premier ($\mathbf{Q}$ ou $\mathbf{F}_p$), les autres s'en déduisant par
\emph{torsion} par un torseur sous le groupe des automorphismes. La question
est en somme celle-ci : les systèmes réguliers de $n$ sommets forment une
\emph{gerbe} sur le corps premier --- sur un corps algébriquement clos ils sont
tous isomorphes, puisque $Q_n$ permute transitivement les invariants ---, et
cette gerbe a-t-elle une section ? En caractéristique $p \mid n$ (donc $n = p$
ou $2p$), existe-t-il un tel système sur $\mathbf{F}_p$ ?\footnote{La page
précise qu'elle considère les structures non nécessairement scindées, sur un
anneau, où $S$ est un sous-schéma fini du fibré affine $E$ ; la parenthèse est
en partie illisible. Le mot « gerbe » est sur la page (page 46), au sens de
Giraud.} Le mieux qu'on puisse espérer :
\begin{enumerate}
  \item[a)] trouver un système régulier de $n$ sommets sur $A =
    \mathbf{Z}[\frac1n]$ ;
  \item[b)] si $n$ est premier impair ou $n = 4$, en trouver un sur $A =
    \mathbf{Z}$ ;
  \item[c)] si $n = 2p$, $p$ premier impair, en trouver un sur $A =
    \mathbf{Z}[\frac12]$\footnote{La page restreint a) à « $n$ premier $\neq
    2$ ou de la forme $2p$ » ; la solution qu'elle donne vaut pour tout $n
    \geq 3$, ce qu'elle dit aussitôt.}.
\end{enumerate}
La solution de a) est « évidente » : dans le plan $\mathbf{E}^2_A$, on prend
\[
  Z_0 = \mu_{n,A} \subset \mathrm{SL}(2)_A, \quad t \mapsto
  \begin{pmatrix} t & 0 \\ 0 & t^{-1} \end{pmatrix}, \qquad
  \sigma_0 = \begin{pmatrix} 0 & 1 \\ 1 & 0 \end{pmatrix}, \qquad
  s_0 = (1, 1), \qquad S_0 = Z_0 \cdot s_0 ,
\]
qui n'est autre que le modèle $s_i = (\zeta^i, \zeta^{-i})$ de la page 8, écrit
sans choisir $\zeta$. Son schéma en groupes d'automorphismes est $G =
\mu_{n,A} \rtimes \{\pm 1\}_A$, une forme tordue du groupe diédral constant, et
la catégorie fibrée des systèmes réguliers de $n$ sommets s'identifie à celle
des $G$-torseurs\footnote{La page note ici $\sigma_0$ l'échange des
coordonnées, qui est notre $\sigma_1$ ; on garde sa lettre dans la formule,
ce morceau ne servant pas ailleurs. L'indice de $\sigma$ dans
$\mathbf{E}^2(A)^{\sigma}$ manque sur la page.}. En caractéristique $p$, on
perd ce modèle, où « les choses sont spéciales » : sur $\mathbf{F}_p$ il y a
bien un système régulier scindé --- les $p$-gones paraboliques et les
$2p$-gones en bande des pages 94 et 95 sont définis sur $\mathbf{F}_p$ --- mais
la question sur $\mathbf{Z}$ (cas b)) ou $\mathbf{Z}[\frac12]$ (cas c)) reste
ouverte. La « Construction » commencée au bas de la page 48 ($T =
\mathrm{Spec}(A)$, $E = \mathbf{E}^2_T$, $s_0$) s'arrête après trois
symboles\footnote{Au-dessus, quelques lignes encadrées et barrées
envisageaient de regarder le cas a) au-dessus de $\mathrm{Spec}(A[\frac1p])$ et
de prolonger. Les pages 114 à 116 reprendront la question sous la forme des
« anneaux modulaires » $k'_n$, et trouveront $k'_4 = \mathbf{Z}$.}.

\pagerange{49}{80}
\section*{II. « Cours Polygones réguliers II » (pages 49 à 80)}

\pagerange{50}{53}
\subsection*{Plan, rappels, et le choix de ne pas se donner de métrique (pages 50 à 52)}

Le plan de la leçon : 1. automorphismes et repères des polygones
combinatoires, le polygone type et les groupes $\mathbf{D}_n$,
$\mathbf{D}_\infty = \mathbf{Z}/2 * \mathbf{Z}/2$ ; 2. espaces affines et
groupe affine $\mathrm{Aff}(E) \simeq \mathrm{GL}(n) \ltimes k^n$ ; 3. métrique
et formes quadratiques, non dégénérées ou non.

Rappel de la leçon précédente : un \emph{polygone géométrique} est un couple
$(S, A)$, $S \subset E$, $A \subset \mathrm{Dr}(E)$, qui est un polygone
combinatoire pour la relation d'incidence ; une \emph{réalisation
géométrique} d'un polygone combinatoire $\Pi = (S, A, R)$ est un couple
$\varphi_S : S \to E$, $\varphi_A : A \to \mathrm{Dr}(E)$ compatible avec
l'incidence. Si deux sommets adjacents ont toujours des images distinctes,
$\varphi_A$ est déterminée par $\varphi_S$, qui est arbitraire. La réalisation
est \emph{fidèle} si $\varphi_S(s) \in \varphi_A(a)$ entraîne que $s$ est
incident à $a$ ; alors $\varphi_S$ et $\varphi_A$ sont injectives et l'image
est un vrai polygone de $E$, dit $\Pi$-épinglé\footnote{C'est le polygone
« propre » de la page 3 ; la page 50 avait d'abord écrit « non dégénérée ».}.

\textbf{Le choix de la définition.} Sur $\mathbf{R}$, avec une métrique, on
avait proposé : le groupe des déplacements qui laissent le polygone invariant
est transitif sur les sommets ; il y a alors un point fixe, le centre de
gravité, les déplacements sont des rotations, et l'on retrouve les $\zeta^i
s_0$ ; $S$ et $A$ sont des torseurs sous ce groupe. Le groupe de \emph{toutes}
les isométries qui conservent $P$ a une transitivité plus forte : il est
(simplement) transitif sur les \emph{repères}. Le cours fait le choix
inverse :
\begin{quote}
Nous préférons ne pas nous donner de métrique d'avance, et prendre les
automorphismes affines de $E$ qui conservent $P$, et \emph{exiger} qu'ils
soient transitifs sur l'ensemble des repères. On récupère la métrique à peu
près de façon unique modulo homothétie --- mais parfois elle sera
dégénérée ! Si on exigeait dès le début une métrique non dégénérée, on
perdrait des cas fort intéressants du point de vue arithmétique --- en fait,
les seuls cas vraiment nouveaux ! Pratiquement, ce nouveau point de vue
consiste à regarder \emph{tous} les triangles comme polygones réguliers, ainsi
que \emph{tous} les parallélogrammes\footnote{Citation des pages 51 et 52,
abréviations développées ; « fort intéressants » est lu avec doute. La
métrique dégénérée est celle des pages 97 et 98 ($\alpha = \pm 2$).}.
\end{quote}
La « définition en forme » annoncée --- polygone régulier, puis réalisation
géométrique régulière d'un polygone combinatoire --- s'arrête à ses deux
titres ; elle est donnée au fil des pages 54 à 80.

\pagerange{54}{61}
\subsection*{Polygones combinatoires, contours et drapeaux (pages 54 à 61)}

Un \emph{polygone combinatoire} est un triple $(S, A, R)$ : un ensemble de
sommets, un ensemble d'arêtes, et une relation d'incidence $R \subset S
\times A$ --- les \emph{repères}, ou drapeaux --- telle que 1) $R \to A$ soit
de degré $2$ (deux sommets par arête), 2) $R \to S$ de degré $2$ (deux arêtes
par sommet), 3) le tout soit connexe. Sans 3) on parle de \emph{contour} ;
sans 2), de graphe ; et si l'on se donne $R \to S$, $R \to A$ sans supposer $R
\hookrightarrow S \times A$, de complexe $1$-cellulaire\footnote{Même matière,
et mêmes mots (« contour »), au dossier 75 (pages 4 à 7). Le langage des
drapeaux est aussi celui du dossier 86.}.

Sur l'ensemble $R$ des drapeaux d'un contour opèrent deux involutions sans
point fixe : $\sigma_1$ échange les deux drapeaux de même sommet ($R/\sigma_1
\simeq S$), $\sigma_0$ les deux drapeaux de même arête ($R/\sigma_0 \simeq
A$). Elles engendrent une action du groupe
\[
  \mathfrak{G}_1 = \langle \sigma_0, \sigma_1 \mid \sigma_0^2 = \sigma_1^2 = 1
  \rangle \simeq \mathbf{Z}/2 * \mathbf{Z}/2 = \mathbf{D}_\infty .
\]

\textbf{Proposition} (page 55). Le foncteur $(S, A, R) \mapsto R$ est une
équivalence de la catégorie des contours sur celle des
$\mathfrak{G}_1$-ensembles « admissibles », ceux où $\sigma_0$ et $\sigma_1$
opèrent sans point fixe et où $\sigma_0 x \neq \sigma_1 x$ pour tout $x$ ---
« c'est une tautologie »\footnote{La dernière condition exprime que $R \to S
\times A$ est injectif ; elle exclut les orbites à deux éléments $\{x,
\sigma_0 x = \sigma_1 x\}$, c'est-à-dire, dit la marge, les « monogones »,
polygones à un côté. La page écrit « dans $E$ » pour « dans $R$ ».}.

Avec $\sigma = \sigma_1$ et $u = \sigma_1 \sigma_0$, $\mathfrak{G}_1 =
\langle \sigma, u \mid \sigma^2 = 1,\ \sigma u \sigma^{-1} = u^{-1} \rangle =
\mathbf{Z} \rtimes \mathbf{Z}/2$, le groupe diédral infini, et les trois
conditions d'admissibilité se résument en
\[
  \forall x \in R, \qquad \operatorname{card}\{x, \sigma x, u x\} = 3
\]
\footnote{Ici $u = \sigma_1\sigma_0$, l'inverse de notre rotation ; à cette
page seulement, cela ne change rien. $\sigma_0 = u^{-1}\sigma = \sigma u$.}.

\textbf{Corollaire} (page 56). Les polygones combinatoires sont les espaces
homogènes de $\mathfrak{G}_1 = \mathbf{D}_\infty$ dont aucun stabilisateur ne
contient $\sigma_0$, $\sigma_1$ ni $u$. Ces stabilisateurs sont exactement les
sous-groupes $n\mathbf{Z} = \langle u^n \rangle$ des rotations, $2 \leq n \leq
\infty$ ($\infty\mathbf{Z} = 0$) : ils sont distingués, de quotient le groupe
diédral $\mathbf{D}_n$. \emph{Démonstration.} La condition est évidemment
suffisante. Si un stabilisateur $H$ contenait un élément $u^m \sigma$ hors des
rotations, comme $u^m \sigma u^{-m} = u^{2m}\sigma$, cet élément serait
conjugué de $\sigma$ ($m$ pair) ou de $u\sigma$ ($m$ impair), donc un
conjugué de $H$, stabilisateur d'un autre point, contiendrait $\sigma_1$ ou
$\sigma_0$. Donc $H \subset \langle u \rangle$, et $u \notin H$.

Soit $\Pi_n$ ($2 \leq n \leq \infty$) le polygone combinatoire défini par
$\mathbf{D}_n$ comme espace homogène de $\mathfrak{G}_1$ : ses sommets sont
$S_n = \langle \sigma_1 \rangle \backslash \mathbf{D}_n$, ses arêtes $A_n =
\langle \sigma_0 \rangle \backslash \mathbf{D}_n$, chacun en bijection avec
$\mathbf{Z}/n$, et $s_i$ est incident à $a_i$ et à $a_{i-1}$\footnote{La
page 57 écrit $S_n = \langle u\sigma \rangle \backslash \mathbf{D}_n$, « $u\sigma
= \sigma_0$ », et $A_n = \langle \sigma \rangle \backslash \mathbf{D}_n$,
« $\sigma = \sigma_1$ » : c'est l'inverse de $S = R/\sigma_1$, $A =
R/\sigma_0$ (page 54), et $u\sigma$ n'est pas $\sigma_0$ avec la convention
$u = \sigma_1\sigma_0$ de la page 55 (il l'est avec celle des pages 77 et
suivantes). Comme $\sigma_0$ et $\sigma_1$ jouent des rôles symétriques dans
$\mathbf{D}_n$, le polygone obtenu est le même aux noms près ; on rétablit la
convention de la page 54. Le croquis montre $a_0$ de $s_0$ à $s_1$.}.

\textbf{Corollaire 2.} Si $\Pi$, $\Pi'$ sont isomorphes et $r \in
\mathrm{Rep}(\Pi)$, $\mathrm{Isom}(\Pi, \Pi') \to \mathrm{Rep}(\Pi')$,
$\varphi \mapsto \varphi(r)$, est bijective --- cas particulier d'un énoncé sur
les espaces homogènes « distingués » (de stabilisateur distingué).

\textbf{Corollaire 3.} Un $n$-gone combinatoire muni d'un repère est défini à
isomorphisme unique près, et canoniquement isomorphe à $\Pi_n$.

\textbf{Corollaire 4.} $\mathrm{Aut}(\Pi_n) \simeq \mathbf{D}_n$
canoniquement (par multiplication à droite) ; $\mathrm{Aut}(\Pi) \simeq
\mathbf{D}_n$, non canoniquement, pour tout $n$-gone $\Pi$.

\textbf{NB.} La restriction $\sigma_0 x \neq \sigma_1 x$ exclut les orbites
$\mathbf{D}_1$, les polygones à un côté ; on pourrait les inclure en
n'exigeant pas $R \hookrightarrow S \times A$. La condition analogue à celle
des polyèdres combinatoires reviendrait à exiger $n \geq 3$, et l'on ne
regardera désormais que $n \geq 3$\footnote{Le bas de la page 58 et le haut de
la page 59 sont en grande partie illisibles ; on n'en retient que la
conclusion, sûre.}.

\textbf{Orientations} (§ 6, pages 59 à 61). Une partie $R^{+} \subset R$ est
une \emph{orientation} si $\sigma_0 R^{+} = \sigma_1 R^{+} = R \smallsetminus
R^{+}$, ou, de façon équivalente, si $R^{+} \to S$ et $R^{+} \to A$ sont
bijectives : $R^{+}$ est le graphe d'une bijection $S \xrightarrow{\sim} A$.
Les orientations d'une somme de contours sont le produit des orientations des
composantes ; un polygone combinatoire a exactement deux orientations,
opposées, qui sont les deux orbites des rotations $\mathfrak{G}_1^{+} =
\langle u \rangle$ dans $R$ (dans $R \simeq \mathbf{D}_n$, les deux classes
$\mathbf{D}_n^{+}$ et $\mathbf{D}_n \smallsetminus \mathbf{D}_n^{+}$).

\textbf{Proposition} (7., pages 60 et 61). Soit $\Pi = (S, A, R)$ un polygone
combinatoire, $G = \mathrm{Aut}(\Pi)$. L'action de $G$ sur les deux
orientations donne une suite exacte
\[
  1 \longrightarrow G^{+} \longrightarrow G \longrightarrow \{\pm 1\}
  \longrightarrow 1 ,
\]
scindée par n'importe quel élément de $G^{-} = G \smallsetminus G^{+}$, qui
sont tous d'ordre $2$, et
\begin{enumerate}
  \item[a)] $G^{+}$ est cyclique d'ordre $n$, et $G^{-}$ opère sur lui par
    inversion ;
  \item[b)] $R$ est un torseur sous $G$, $S$ et $A$ des torseurs sous $G^{+}$ ;
  \item[c)] pour toute orientation $\omega$, il existe un unique $u_\omega \in
    G^{+}$ qui coïncide avec $u$ sur $R^{+}_\omega$ ; $u_{-\omega} =
    u_\omega^{-1}$, $u_\omega$ engendre $G^{+}$, et si $n \geq 3$,
    $\omega \mapsto u_\omega$ est une injection de $\mathrm{Or}(\Pi)$ dans
    l'ensemble des générateurs de $G^{+}$ ;
  \item[d)] le foncteur $\Pi \mapsto (G^{+}, \{u_\omega, u_{-\omega}\}, S)$ est
    une équivalence de la catégorie des $n$-gones ($3 \leq n \leq \infty$)
    sur celle des triples formés d'un groupe cyclique d'ordre $n$ (monogène
    infini si $n = \infty$), d'une paire de générateurs inverses l'un de
    l'autre, et d'un torseur sous ce groupe\footnote{Le premier terme du
    triple se lit « $\mathbf{Z}\,G^{+}$ », peut-être $\mathbf{Z} \simeq
    G^{+}$ pour $n = \infty$ ; la note marginale ajoute « ou ce qui revient au
    même $\mathbf{Z}$-torseur ». C'est la description par $(G^0, \{g,
    g^{-1}\}, S)$ du dossier 75 (page 6).}.
\end{enumerate}
\emph{Exercice} : exprimer $A$, et $R$ comme partie de $S \times A$, en termes
de ces données\footnote{La page écrit « partie de $S \times R$ », la lecture
du dernier facteur étant sûre ; il faut $S \times A$. L'exercice n'est pas
fait. La réponse attendue : $A$ est l'ensemble des paires $\{s, u_\omega s\}$,
et $R$ celui des couples $(s, a)$ avec $s \in a$.}.

\pagerange{62}{73}
\subsection*{Le cas d'un corps algébriquement clos, $p \nmid n$ (pages 62 à 73)}

Soit $k$ algébriquement clos de caractéristique $p$ ne divisant pas $n$, et
$\Pi$ un $n$-gone régulier ($3 \leq n < \infty$) d'un plan $E$. Le barycentre
$o$ de $S$ --- qui a un sens parce que $n$ est inversible dans $k$ --- est
fixé par $G = \mathrm{Aut}(\Pi)$ ; prenant $o$ pour origine, $E$ devient un
plan vectoriel où $G$ opère linéairement. Soit $u$ l'un des deux générateurs
privilégiés de $G^{+}$, $s_0$ un sommet, $\sigma_1$ l'automorphisme $\neq 1$
qui fixe $s_0$ : $\sigma_1^2 = 1$, $\sigma_1 u \sigma_1 = u^{-1}$, $u^n = 1$,
$n$ minimal.

Soient $\zeta$, $\zeta'$ les valeurs propres de $u$. Si $\zeta = \zeta'$, on
écrit $u = \zeta(\mathrm{id} + v)$ avec $v$ nilpotent ; $u^n = \zeta^n
(\mathrm{id} + n v) = \mathrm{id}$ donne $nv = 0$, donc $v = 0$ ($n$
inversible), $u$ est une homothétie et les $s_i$ seraient alignés sur la
droite $o s_0$ : absurde. Donc $\zeta \neq \zeta'$, d'où deux droites
propres $D_\zeta$, $D_{\zeta'}$. Comme $u$ est conjugué de $u^{-1}$, l'ensemble
$\{\zeta, \zeta'\}$ est stable par inversion, et $\zeta' =
\zeta^{-1}$\footnote{La stabilité par inversion laisse une autre
possibilité, $\{\zeta, \zeta'\} = \{1, -1\}$, que la page ne mentionne pas ;
elle donnerait $u^2 = 1$, exclu puisque $u$ est d'ordre $n \geq 3$.}.
$\sigma_1$ échange $D_\zeta$ et $D_{\zeta^{-1}}$ ; écrivant $s_0 = e_1 + e_2$
($e_1 \in D_\zeta$, $e_2 \in D_{\zeta^{-1}}$), $\sigma_1 s_0 = s_0$ donne
$\sigma_1 e_1 = e_2$, et $e_1, e_2 \neq 0$ (sinon $S$ serait sur une droite
propre). Dans la base $(e_1, e_2)$, attachée à l'épinglage,
\[
  u = \begin{pmatrix} \zeta & 0 \\ 0 & \zeta^{-1} \end{pmatrix}, \qquad
  \sigma_1 = \begin{pmatrix} 0 & 1 \\ 1 & 0 \end{pmatrix}, \qquad
  s_i = \begin{pmatrix} \zeta^i \\ \zeta^{-i} \end{pmatrix} .
\]
On retrouve le cas typique de la page 8, sans avoir supposé la caractéristique
nulle ni $\zeta$ d'ordre fini. L'invariant $\alpha = \zeta + \zeta^{-1} =
\operatorname{Tr} u = \operatorname{Tr} u^{-1}$, ou la paire non ordonnée
$\{\zeta, \zeta^{-1}\}$, est défini intrinsèquement par $\Pi$.

\textbf{Proposition} (pages 64 et 65). Soit $k$ algébriquement clos de
caractéristique $p$ ne divisant pas $n$, $3 \leq n < \infty$.
\begin{enumerate}
  \item[a)] Tout $n$-gone régulier épinglé sur $k$ est isomorphe à l'un des
    polygones épinglés $\Pi_{\zeta, k}$ de $k^2$ ci-dessus, $\zeta$ racine
    primitive $n$-ième de $1$ définie au passage près à $\zeta^{-1}$.
  \item[b)] D'où des bijections entre les classes d'isomorphie de $n$-gones
    réguliers (épinglés ou non), les paires $\{\zeta, \zeta^{-1}\}$ de racines
    primitives $n$-ièmes de $1$, et les $\alpha \in k$ tels que les deux
    racines de $\zeta^2 - \alpha \zeta + 1 = 0$ soient des racines primitives
    $n$-ièmes de $1$ --- les « invariants numériques » des polygones
    réguliers.
  \item[c)] Dans $\Pi_{\zeta, k}$, l'origine est le barycentre de $S$ et
    l'unique point fixe de $G$, et de $G^{+}$ ; les axes de coordonnées sont
    les droites propres des générateurs privilégiés de $G^{+}$, de valeurs
    propres $\zeta$ et $\zeta^{-1}$\footnote{La fin de la page 64 est
    surchargée d'ajouts ; on ne lit sûrement que « l'unique point fixe sous
    $G$ », « $G^{+}$ » et « $\{\zeta, \zeta^{-1}\}$ » ; on complète c) par ce
    que la page 63 établit.}.
\end{enumerate}

\textbf{Formes invariantes.} Les fonctions polynomiales invariantes par $G$
sont sommes de leurs composantes homogènes. Comme $u^{*}(x^a y^b) =
\zeta^{a-b} x^a y^b$, une forme $\sum c_{ab}\, x^a y^b$ est invariante par
$u$ si et seulement si $c_{ab} = 0$ pour $a \not\equiv b \pmod n$ ;
l'invariance par $\sigma_1$ est la symétrie en $x, y$. Les invariants ont donc
pour base les $(xy)^a$ et les $(xy)^a (x^{in} + y^{in})$, $a \in \mathbf{N}$,
$i \geq 1$\footnote{La page donne pour base $1$ et les $(xy)^a(x^{in} +
y^{in})$, $a \in \mathbf{N}$, $i \in \mathbf{N}^{*}$, oubliant les $(xy)^a$,
$a \geq 1$ (le cas $i = 0$, qui donnerait $2(xy)^a$).}. Comme $X^i + Y^i =
S'_i(X + Y, XY)$ pour un polynôme universel $S'_i$ à coefficients entiers
(fonctions symétriques), on obtient
\[
  k[x, y]^{G} = k[\,xy,\ x^n + y^n\,] .
\]
\emph{NB.} Le calcul marche pour tout $n \geq 1$ et sur tout anneau de base,
« le meilleur étant $\mathbf{Z}$ » : sous la forme qu'on dirait aujourd'hui,
l'anneau des invariants du schéma en groupes $\mu_n \rtimes \mathbf{Z}/2$
opérant sur $\mathbf{A}^2_{\mathbf{Z}}$ (comme à la page 47) est
$\mathbf{Z}[xy, x^n + y^n]$\footnote{La page dit simplement que les calculs
« marchent pour tout $n \in \mathbf{N}^{*}$ \ldots{} et sur tout anneau de
base ». Pour un $\zeta$ donné dans un anneau, il faut que les $\zeta^d - 1$
($0 < d < n$) ne soient pas diviseurs de zéro ; la formulation par le schéma
en groupes évite cette réserve. Pour $n$ infini, l'anneau des invariants est
$k[xy]$ (page 67).}.

\textbf{Corollaire} (page 67). Pour $n \geq 3$, les fonctions polynomiales de
degré $\leq 2$ invariantes par $G$ sont les $a\,xy + b$. À un scalaire près,
il y a donc une unique forme quadratique invariante, $q_0 = xy$, normalisée
par $q_0(s) = 1$ pour $s \in S$, et une unique conique \emph{invariante}
circonscrite, $xy - 1 = 0$. Pour $n \geq 5$, c'est même l'unique conique
circonscrite, cinq points d'une conique non dégénérée la
déterminant\footnote{C'est ce que dit la marge, en partie illisible : « Ne
marche que si $n \geq 5$ \ldots{} Ici on a une conique circonscrite canonique
--- d'équation invariante par $G$ ».}.

Pour un $n$-gone régulier non épinglé, il y a donc une unique forme
quadratique invariante $q_0$ normalisée par $q_0(s) = 1$, et une unique
conique circonscrite invariante $C : q_0(x) = 1$. La forme $q_0$, vue sur
l'espace des translations, est non dégénérée : elle fait de $E$ un « plan
euclidien ». Le carré de la longueur d'un côté est
\[
  q_0(s_1 - s_0) = (\zeta - 1)(\zeta^{-1} - 1) = 2 - \alpha ,
\]
non nul puisque $\zeta \neq 1$. On pourrait normaliser plutôt par $q_0(a) =
1$ ; « dans le cas général, c'est cette normalisation-là qui sera la
meilleure » --- c'est celle de la forme $f_0$ de la page 97.

\textbf{Remarque} (pages 69 et 70). Si l'on avait seulement exigé que le
groupe des automorphismes affines soit transitif sur les sommets, on aurait
trouvé beaucoup plus : un $u$ d'ordre $n$ sans $\sigma$, c'est-à-dire deux
valeurs propres $\zeta$, $\zeta'$ racines de l'unité d'ordres $d$, $d'$ avec
$n = \operatorname{ppcm}(d, d')$, et $S$ l'orbite d'un point hors des deux
droites propres. Exiger de plus une métrique non dégénérée invariante dont
les automorphismes soient des \emph{rotations} « sauve la mise » : une
rotation est unimodulaire, ses valeurs propres sont $\{\zeta, \zeta^{-1}\}$,
et l'on retombe sur les cas précédents\footnote{La page attribue la
suggestion à un collègue dont le nom est lu avec doute ; on ne le restitue
pas. Le paragraphe barré qui précède (« Commentaires philosophiques », pages
68 et 69) décrivait la famille, paramétrée par $\alpha$, des représentations
affines de dimension $2$ de $\mathbf{D}_n$ ; il est abandonné.}. Il ne sait pas
si l'on simplifie encore en exigeant seulement que les éléments de $G^{+}$
respectent la métrique.

\textbf{Semi-invariants} (pages 71 à 73). Les monômes $x^i y^j$ sont vecteurs
propres de $u$, de valeur propre $\zeta^{i-j}$. L'algèbre des invariants de
$G^{+}$ est
\[
  k[x, y]^{G^{+}} = k[W, U, V], \qquad W = xy,\ U = x^n,\ V = y^n,\qquad UV =
  W^n ,
\]
et, pour $0 < d < n$, la composante de valeur propre $\zeta^d$ est le module
engendré par $X_d = x^d$ et $Y_d = y^{n-d}$, avec la relation $V X_d = W^d
Y_d$. En degré $\leq 2$, les valeurs propres de $1, x, y, x^2, xy, y^2$ sont
$1, \zeta, \zeta^{-1}, \zeta^2, 1, \zeta^{-2}$.
\begin{itemize}
  \item Si $n \geq 5$, ces valeurs $\zeta^{-2}, \ldots, \zeta^2$ sont
    distinctes ; les semi-invariants de $G^{+}$ de degré $\leq 2$ sont $a +
    bxy$, $ax$, $ay$, $ax^2$, $ay^2$, et comme $\sigma_1$ échange $x$ et $y$,
    seuls les $a + bxy$ sont semi-invariants (et même invariants) sous $G$.
  \item Si $n = 3$ ($\zeta^2 = \zeta^{-1}$), les semi-invariants de $G^{+}$
    sont $a + bxy$, $ax + by^2$ (valeur propre $\zeta$), $ay + bx^2$
    ($\zeta^{-1}$) ; sous $G$, encore les seuls $a + bxy$.
  \item Si $n = 4$ ($\zeta^2 = \zeta^{-2} = -1$), il s'ajoute $ax^2 + by^2$
    (valeur propre $-1$), et parmi eux $x^2 + y^2$ et $x^2 - y^2$ sont
    semi-invariants sous $G$. Les coniques semi-invariantes circonscrites au
    carré sont $xy - 1 = 0$ et $x^2 - y^2 = 0$, la paire de ses
    diagonales\footnote{La page dit « dans chacune [des familles], il y a une
    conique circonscrite à $S$, savoir l'habituelle $xy - 1 = 0$, et de plus
    $x^2 - y^2 = 0$ » ; que celle-ci soit la paire des diagonales $y = \pm x$
    du carré $(\pm 1, \pm 1)$, $(\pm i, \mp i)$ est de nous.}.
\end{itemize}
Les seules coniques invariantes par $G$ sont donc les $a + bxy = 0$, et la
seule qui soit circonscrite est $xy - 1 = 0$ ; le choix des coordonnées
épinglées universelles montre, ajoute-t-il, que c'est la seule conique
circonscrite\footnote{Pour $n \geq 5$ ; pour $n = 3, 4$, il y a d'autres
coniques circonscrites, non invariantes. La page 72 écrit la condition de
passage par les sommets « $a + b = 0$ », avec la normalisation $xy = 1$ sur
$S$.}.

\pagerange{74}{80}
\subsection*{Réaliser $\Pi_\infty$ : programme et reformulation (pages 74 à 80)}

\textbf{Programme.} Soit $\Pi$ un polygone combinatoire ; on étudie ses
réalisations géométriques régulières dans des plans affines variables. Deux
réalisations $\Pi$-isomorphes le sont par un isomorphisme \emph{unique} ; dans
chaque classe, la réalisation $(E_\alpha, \varphi_{\alpha,S},
\varphi_{\alpha,A})$, et en particulier son plan, est donc définie à
isomorphisme unique près. On indexera les classes par un invariant numérique
$\alpha$, et il faudra écrire explicitement $E_\alpha$, $\varphi_{\alpha,S}$,
$\varphi_{\alpha,A}$ en fonction de $\Pi$ et de $\alpha$. Épinglé par un
repère, $\Pi$ devient canoniquement isomorphe à un $\Pi_n$, $3 \leq n \leq
\infty$, et tous les $\Pi_n$ sont des quotients canoniques de $\Pi_\infty$ (et
$\Pi_n$ de $\Pi_m$ si $n \mid m$), compatibles aux repères : la catégorie
(discrète) des réalisations régulières de $\Pi_n$ est une sous-catégorie pleine
de celle de $\Pi_\infty$. D'où trois étapes :
\begin{enumerate}
  \item[a)] étudier les réalisations régulières de $\Pi_\infty$ sur un corps,
    plus généralement un anneau ; on verra qu'elles sont caractérisées par un
    invariant $\alpha \in k$, arbitraire ;
  \item[b)] expliciter les conditions sur $\alpha$ pour que la réalisation se
    factorise par $\Pi_n$, et de façon fidèle --- d'où une classification des
    $n$-gones réguliers sur $k$ par un invariant $\alpha$ soumis à des
    conditions ;
  \item[c)] en déduire les propriétés essentielles et la classification des
    $n$-gones réguliers, sans épinglage ; « cette étape est essentiellement
    triviale ».
\end{enumerate}
L'étape a) occupe les pages 76 à 100, l'étape b) les pages 101 à 120.

\textbf{Reformulation.} Avec $G = \mathrm{Aut}(\Pi_\infty) \simeq
\mathbf{D}_\infty$, une réalisation régulière de $\Pi_\infty$ dans $E$ est
donnée par
\[
  (2) \quad \sigma_0, \sigma_1 \in \mathrm{Aff}(E), \ \sigma_0^2 = \sigma_1^2
  = \mathrm{id}_E ; \qquad
  (3) \quad r = (s_0, a_0) \in E \times \mathrm{Dr}(E),\ s_0 \in a_0 ,
\]
le drapeau image du repère standard $r_\infty$, d'où, avec $u = \sigma_0
\sigma_1$, $\varphi_S(\underline{s}_n) = u^n s_0$ et
$\varphi_A(\underline{a}_n) = u^n a_0$. L'application $\varphi_{\mathrm{rep}}
: \mathrm{Rep}(\Pi_\infty) \to \mathrm{Drap}(E)$, $g\, r_\infty \mapsto
\varphi_G(g)\, r$, vers l'ensemble des drapeaux affines (point, droite
incidente) de $E$, doit passer au quotient en $\varphi_S$ et $\varphi_A$ :
\begin{tikzcd}
\mathrm{Rep}(\Pi_\infty) \arrow[r, "\varphi_{\mathrm{rep}}"] \arrow[d] & \mathrm{Drap}(E) \arrow[d] \\
S(\Pi_\infty) \times A(\Pi_\infty) \arrow[r, "\varphi_S \times \varphi_A"'] & E \times \mathrm{Dr}(E)
\end{tikzcd}
ce qui équivaut à
\[
  (8) \qquad \sigma_1 s_0 = s_0, \qquad \sigma_0 a_0 = a_0
\]
\footnote{La page 78 dessine deux carrés superposés, $\mathrm{pr}_1$ vers
$S(\Pi_\infty) \to E$ et $\mathrm{pr}_2$ vers $A(\Pi_\infty) \to
\mathrm{Dr}(E)$ ; on les réunit en un seul, par le produit. « $\mathrm{Drap.aff}(E)$ »
est une lecture probable. Les sommets et arêtes combinatoires sont soulignés
sur la page ; on garde le soulignement ici seulement.}. Posons $s_1 = \sigma_0
s_0 = u s_0$ et $s_{-1} = \sigma_1 s_1 = u^{-1} s_0$.

\textbf{Lemme} (page 79). Étant donnés (2) et (3) satisfaisant (8), les
conditions suivantes sont équivalentes :
\begin{enumerate}
  \item[(a)] les données définissent une réalisation de $\Pi_\infty$,
    c'est-à-dire que les $s_n = u^n s_0$ engendrent affinement $E$ ;
  \item[(b)] $s_0$, $s_1$, $s_{-1}$ sont affinement indépendants ;
  \item[(c)] $\varphi_S$ et $\varphi_A$ ne sont pas constantes\footnote{La
    page écrit en (a) que les $s_n$ sont « affinement indépendants », ce qui
    pour une suite infinie de points d'un plan ne peut s'entendre que comme
    « engendrent affinement $E$ ».}.
\end{enumerate}
\emph{Démonstration.} (b) $\Rightarrow$ (a) $\Rightarrow$ (c) sont
immédiats. Supposons (c). 1°) $s_1 \neq s_0$, sinon $s_0$ serait fixé par
$\sigma_0$ et $\sigma_1$, donc par $G$, et $\varphi_S$ serait constante.
2°) Alors $a_0 = \mathrm{dr}(s_0, s_1)$, car $s_0 \in a_0$ et $\sigma_0 a_0 =
a_0$ donne $s_1 = \sigma_0 s_0 \in a_0$ (9) ; et $\sigma_1 a_0 =
\mathrm{dr}(s_0, s_{-1})$. Si $s_{-1}$ était sur $a_0$, on aurait $\sigma_1
a_0 = a_0$, $a_0$ serait fixée par $G$ et $\varphi_A$ serait
constante\footnote{La page écrit « $\sigma_0 a_0 = \mathrm{dr}(\sigma_0 s_0 =
s_1, \sigma_1 s_1 = s_{-1})$ » et « si on avait $s_{-1} \in a_1$ » : il faut
$\sigma_1 a_0$ et $a_0$, comme la suite l'exige.}.

\textbf{Proposition} (page 80). Les réalisations géométriques régulières de
$\Pi_\infty$ dans $E$ correspondent bijectivement aux triples $(\sigma_0,
\sigma_1, s_0)$, $\sigma_i \in \mathrm{Aff}(E)$, $s_0 \in E$, tels que
\begin{enumerate}
  \item[a)] $\sigma_0^2 = \sigma_1^2 = \mathrm{id}$ ;
  \item[b)] $\sigma_1 s_0 = s_0$ ;
  \item[c)] $s_0$, $s_1 = \sigma_0 s_0$ et $s_{-1} = \sigma_1 s_1$ forment un
    repère affine de $E$.
\end{enumerate}
La représentation associée envoie les générateurs $\underline{\sigma}_i$ de
$\mathrm{Aut}(\Pi_\infty)$ sur $\sigma_i$, et $\varphi_S(\underline{s}_n) = u^n
s_0$, $\varphi_A(\underline{a}_n) = \mathrm{dr}(s_n, s_{n+1}) = u^n
\mathrm{dr}(s_0, s_1)$, $n \in \mathbf{Z}$. C'est la proposition 5 du
chapitre I, à ceci près que la non-colinéarité porte sur $s_{-1}, s_0, s_1$, ce
qui revient au même.

\pagerange{81}{100}
\section*{III. Les coordonnées universelles (pages 81 à 100)}

\pagerange{81}{85}
\subsection*{Les paramètres $\alpha$, $\beta$ ; cas propre et cas impropre (pages 81 à 85)}

Le repère $(s_0, s_1, s_{-1})$ identifie $E$ à $k^2$ : origine $s_0$, base
$e_1 = s_1 - s_0$, $e_2 = s_{-1} - s_0$, et $(x, y) = s_0 + x e_1 + y e_2$.
Ainsi $s_0 = (0,0)$, $s_1 = (1, 0)$, $s_{-1} = (0, 1)$\footnote{La page 81
écrit « $s_2 = (0,1)$ » ; par (15), c'est $s_{-1}$.}. L'automorphisme
$\sigma_1$ fixe $s_0$ et échange $s_1$, $s_{-1}$ :
\[
  (21) \qquad \sigma_1(x, y) = (y, x) .
\]
L'automorphisme $\sigma_0$ échange $s_0$ et $s_1$ et envoie $s_{-1}$ sur $s_2 =
u^2 s_0 = \sigma_0 s_{-1}$ ; on \emph{définit} $\alpha, \beta \in k$ par
\[
  s_2 = s_0 + (1 + \alpha)\, e_1 + (1 + \beta)\, e_2 ,
\]
d'où $\sigma_0 e_1 = -e_1$, $\sigma_0 e_2 = \alpha e_1 + (1 + \beta) e_2$, et
\[
  (17), (18) \qquad \sigma_0 = \begin{pmatrix} -1 & \alpha & 1 \\ 0 & 1 + \beta
  & 0 \\ 0 & 0 & 1 \end{pmatrix}, \qquad \sigma_0(x, y) = (1 - x + \alpha y,\
  (1 + \beta) y)
\]
(en coordonnées homogènes $(x, y, 1)$). Reste à écrire $\sigma_0^2 =
\mathrm{id}$ : $\sigma_0^2(x, y) = (x + \alpha\beta\, y,\ (1 + \beta)^2 y)$,
donc
\[
  (22) \qquad \alpha\beta = 0, \qquad \beta(\beta + 2) = 0 .
\]

\textbf{Proposition} (page 83). Les réalisations géométriques régulières de
$\Pi_\infty$ sur $k$, épinglées, sont en bijection avec les couples $(\alpha,
\beta) \in k^2$ satisfaisant (22).

L'exemple $s_i = (\zeta^i, \zeta^{-i})$ (page 82) donne $s_2 - s_{-1} = (1 +
\zeta + \zeta^{-1})(s_1 - s_0)$, la relation de la page 7 ; on a donc envie
d'imposer $\beta = 0$.

\textbf{Proposition} (pages 82 et 84). Pour une réalisation régulière de
$\Pi_\infty$, les conditions suivantes sont équivalentes :
\begin{enumerate}
  \item[(i)] $\beta = 0$ ;
  \item[(ii)] $s_2 - s_{-1} = (1 + \alpha)(s_1 - s_0)$ ; (ii$'$) $s_2 -
    s_{-1} \parallel s_1 - s_0$ ;
  \item[(iii)] pour tout $n \in \mathbf{Z}$, $s_{n+3} - s_n = (1 + \alpha)(s_{n+2}
    - s_{n+1})$ ; (iii$'$) pour tout $n$, $s_{n+3} - s_n \parallel s_{n+2} -
    s_{n+1}$\footnote{La page écrit (iii$'$) mot pour mot comme (iii) ; on
    l'écrit comme l'appelle (ii$'$).} ;
  \item[(iv)] $\det \sigma_{0V} = -1$ ; (iv$'$) $\det u_V = 1$ ; (iv$''$)
    $\operatorname{Tr} \sigma_{0V} = 0$ ;
  \item[(v)] on n'est pas dans le cas \emph{impropre} : $\operatorname{car} k
    \neq 2$, $\alpha = 0$, $\beta = -2$, $\sigma_0(x, y) = (1 - x, -y)$,
    symétrie par rapport au milieu de $[s_0, s_1]$.
\end{enumerate}
\emph{Démonstration.} $s_2 - s_{-1} = (1 + \alpha) e_1 + \beta e_2$, d'où (i)
$\Leftrightarrow$ (ii) $\Leftrightarrow$ (ii$'$) ; (iii) donne (ii) pour $n =
-1$ et s'en déduit en appliquant $u^{n+1}$. Comme
\[
  (24) \qquad \det u_V = -\det \sigma_{0V} = 1 + \beta, \qquad
  \operatorname{Tr} \sigma_{0V} = \beta ,
\]
(iv), (iv$'$), (iv$''$) équivalent à (i) ; et (v) à (i) par (22).

La condition « propre » est donc exactement l'\emph{unimodularité} du
chapitre I.

\textbf{Le cas impropre} (pages 84 et 85). À isomorphisme près il y en a un et
un seul pour $k$ donné ($\operatorname{car} k \neq 2$) :
\[
  (25) \qquad u(x, y) = (1 - y, -x), \qquad u^2(x, y) = (x + 1, y - 1) ,
\]
$u^2$ est la translation de vecteur $(1, -1)$. En caractéristique $0$, $u$ est
d'ordre infini, et l'on obtient un $\infty$-polygone régulier fidèle : les
sommets $s_{2m} = (m, -m)$ et $s_{2m+1} = (m + 1, -m)$, en escalier. En
caractéristique $p > 2$, $u^2$ est d'ordre $p$, $u$ d'ordre $2p$, et la
réalisation impropre se factorise par $\Pi_{2p}$ en une réalisation fidèle :
un $2p$-gone régulier, non unimodulaire\footnote{Les $2p$ points $(m, -m)$,
$(m+1, -m)$, $m \in \mathbf{F}_p$, sont distincts ; c'est le calcul de la page
85, « on vérifie aisément », fait ici.}. « Avec la théorie projective des
polygones réguliers, l'anomalie des polygones réguliers impropres s'explique
complètement », et l'on n'a rien perdu, du point de vue projectif, à
l'exclure, ce qu'on fera désormais\footnote{La théorie projective n'est pas
dans le dossier. Le mot qualifiant « anomalie » est illisible.}.

\textbf{Scholie} (page 85). Une réalisation géométrique régulière propre de
$\Pi_\infty$ est décrite par $\alpha \in k$, qui peut être choisi
arbitrairement. La marge baptise $\alpha$ « invariant numérique ».

\pagerange{86}{89}
\subsection*{Points fixes, centre, valeurs propres (pages 86 à 89)}

Dans le cas propre ($\beta = 0$) :
\[
  (26) \qquad \sigma_0(x, y) = (1 - x + \alpha y,\ y), \quad
  \sigma_1(x, y) = (y, x), \quad u(x, y) = (1 - y + \alpha x,\ x) .
\]
Le lieu fixe de $\sigma_0$ est la droite
\[
  (27) \qquad H_0 : \ 2x - \alpha y - 1 = 0 ,
\]
sauf en caractéristique $2$ avec $\alpha = 0$ (le carré de caractéristique
$2$), où $H_0 = \emptyset$ et $\sigma_0$ est la translation de vecteur $e_1$ ;
en coordonnées homogènes on trouve encore une droite, la droite à l'infini.
Celui de $\sigma_1$ est la droite $H_1 : x = y$. Appelant \emph{réflexion
affine} un automorphisme d'un plan affine dont le lieu fixe est une droite, le
cas propre est donc encore caractérisé par
\begin{enumerate}
  \item[(iv$'''$)] $\sigma_0$ est une réflexion affine, ou une translation
    (dans le seul cas $p = 2$, $\alpha = 0$, où elle reste une réflexion
    projective) ;
  \item[(iv$^{\mathrm{iv}}$)] $\sigma_{0V}$ est une réflexion, $\neq
    \mathrm{id}$\footnote{La page écrit « sauf pour $p = 2$, $\alpha = 2$ »,
    puis, à la ligne (iv$'''$), « $p = 2$, $\alpha = 0$ » : en caractéristique
    $2$ c'est la même valeur. Dans le cas impropre, $\sigma_0$ est la symétrie
    par rapport à un point, et $\sigma_{0V} = -\mathrm{id}$.}.
\end{enumerate}

Les points fixes de $G$, de $G^{+}$ et de $u$ coïncident ; ce sont les
solutions de $x = y$, $x(2 - \alpha) = 1$ :
\[
  (29) \qquad x_0 = y_0 = \frac{1}{2 - \alpha} ,
\]
qui existe, unique, si et seulement si $\alpha \neq 2$. C'est le \emph{centre}
de la réalisation.

La partie linéaire de la rotation est
\[
  (30) \qquad u_V = \begin{pmatrix} \alpha & -1 \\ 1 & 0 \end{pmatrix}, \qquad
  \det u_V = 1, \quad \operatorname{Tr} u_V = \alpha ,
\]
d'équation aux valeurs propres $\zeta^2 - \alpha\zeta + 1 = 0$ ($\zeta \in
\bar k$), de discriminant $\Delta(\alpha) = \alpha^2 - 4 = (\alpha -
2)(\alpha + 2)$ ; ses racines vérifient $\zeta\zeta' = 1$ et $\alpha = \zeta +
\zeta^{-1}$. L'invariant numérique est donc celui de la page 63, défini
maintenant sans condition. Le cas d'une racine double, $\zeta = \zeta' = \pm
1$, c'est-à-dire $\alpha = \pm 2$ ou $\Delta(\alpha) = 0$, n'est plus exclu
($\alpha = 2$ pour $\zeta = 1$, $\alpha = -2$ pour $\zeta = -1$, confondus
si et seulement si $p = 2$) : ce sont exactement les deux valeurs de $\zeta$
qu'il avait fallu écarter pour définir la représentation $\varphi^\zeta$, et
$\alpha = 2$ est le seul cas sans centre à distance finie.

Pour $\alpha \neq \pm 2$, et $\zeta \in k$, on retrouve la représentation
$\varphi^\zeta$ ; sur $k$ algébriquement clos, les seuls cas nouveaux sont donc
$\alpha = \pm 2$. Et il y insiste :
\begin{quote}
même si on décidait (à tort !) d'exclure ces deux cas, les considérant comme
trop « exotiques », ils y reviendraient indirectement, car a) ils sont
valables pour tout corps, et même sur les anneaux de base, pas seulement pour
$k$ algébriquement clos, et b) ils permettent de suivre la variation de
$\varphi(\alpha)$ en termes des paramètres $\alpha$,
intrinsèques\footnote{Page 89, abréviations développées ; trois mots
illisibles après « décidait » sont rendus par « d'exclure », que le sens
impose, et « même sur les » est lu avec doute.}.
\end{quote}

\pagerange{90}{95}
\subsection*{Étude par cas (pages 90 à 95)}

\textbf{I. $\alpha \neq \pm 2$, cas ordinaire.} Étendant au besoin $k$ en
$k(\zeta)$ et rapportant $E$ au centre et aux droites propres $D_\zeta$,
$D_{\zeta^{-1}}$ de $u$, avec $s_0 = (1,1)$, on retrouve le cas type
$\Pi_\zeta$, où $\zeta$ n'est pas forcément racine de l'unité. Il y a une forme
quadratique invariante non dégénérée, unique à un scalaire près. L'ordre de
$u$ est celui de $u_V$, c'est-à-dire celui de $\zeta$ ; quand il est fini,
c'est le nombre de côtés du polygone image, et il est premier à la
caractéristique.

\textbf{II. $\alpha = -2$, $p \neq 2$.} La partie linéaire s'écrit $u_V =
-(\mathrm{id} + v_0)$, $v_0^2 = 0$, $v_0 \neq 0$ ; plus généralement, pour
$\alpha = 2\varepsilon$, $\varepsilon = \pm 1$, $u_V = \varepsilon(\mathrm{id}
+ v_0)$ avec
\[
  v_0(x, y) = (x - \varepsilon y)\,(1, \varepsilon) = a'(x, y)\, a, \qquad
  a' = x - \varepsilon y \in V^{*}, \quad a = (1, \varepsilon) \in V ,
\]
un endomorphisme de rang $1$ et de carré nul\footnote{Le cas II est d'abord
écrit puis barré (page 90), repris page 91, essayé sur une page de brouillon
(page 92) et récrit page 94 après quatre départs biffés ; on suit la version
de la page 94.}. Il y a un centre, $(\frac14, \frac14)$, et dans les
coordonnées centrées $X = x - \frac14$, $Y = y - \frac14$,
\[
  \sigma_0(X, Y) = (-X - 2Y,\ Y), \quad \sigma_1(X, Y) = (Y, X), \quad
  u = -(\mathrm{id} + v_0), \quad v_0 = (X + Y)\,(1, -1) .
\]
Comme $u^{2m} = \mathrm{id} + 2m\, v_0$ et $u^{2m+1} = -(\mathrm{id} +
(2m+1) v_0)$, les sommets sont
\[
  s_{2m} = s_0 + m\, b, \qquad s_{2m+1} = s_1 - m\, b, \qquad b = (-1, 1)
\]
dans les coordonnées de départ ($s_2 = (-1, 1)$, $s_3 = (2, -1)$, \ldots) :
les sommets pairs sur la droite $x + y = 0$, les impairs sur la droite
parallèle $x + y = 1$, parcourues en sens contraires --- un polygone qui
zigzague dans une bande\footnote{La page 95 écrit, en coordonnées centrées,
$s_0 = (\frac14, \frac14)$, $s_{2n} = s_0 + n a$, $s_{2n+1} = (-s_0 - \frac12
a) - na$ avec $a = (1, -1)$ ; les signes sont des pâtés lus d'après le calcul.
En coordonnées centrées $s_0 = (-\frac14, -\frac14)$ (le centre est en
$(\frac14, \frac14)$ dans les coordonnées de départ, où $s_0$ est l'origine),
et la description par deux droites parallèles parcourues en sens contraires
ne dépend pas de ce signe. On la donne dans les coordonnées de départ, où elle
se vérifie directement sur $u(x, y) = (1 - y - 2x, x)$.}. En
caractéristique $0$ c'est un $\infty$-polygone ; en caractéristique $p > 2$,
$u_V^2 = \mathrm{id} + 2v_0$ est une transvection d'ordre $p$, $u_V$ et $u$
sont d'ordre $2p$ : un $2p$-gone régulier (unimodulaire), propre, distinct de
l'impropre de la page 85.

En caractéristique $2$, $\alpha = -2 = 0 = 2$ : c'est le \emph{carré},
$u(x, y) = (1 - y, x)$, $u^2(x, y) = (1 - x, 1 - y)$, $u^4 = \mathrm{id}$,
d'ordre $4 = 2p$ ; il n'a pas de centre.

\textbf{III. $\alpha = 2$, pas de centre.} $u_V = \mathrm{id} + v_0$, $v_0 =
a' \otimes a \neq 0$, $v_0^2 = 0$ : en caractéristique $p$, $u_V$ est d'ordre
$p$. Dans l'enveloppe vectorielle $\hat E$ de $E$ (le plan affine plongé comme
hyperplan $z = 1$ de $k^3$), $u_{\hat E} = \mathrm{id} + w$ avec $w^3 = 0$ et
$w^2 \neq 0$, puisque $u$ n'a pas de point fixe ; donc $u^p = \mathrm{id} +
w^p = \mathrm{id}$ si $p \geq 3$, et $u$ est d'ordre $p$ ; pour $p = 2$,
$u^2 = \mathrm{id} + w^2 \neq \mathrm{id}$ et l'on retrouve le carré, d'ordre
$4$. En caractéristique $0$, $u$ est d'ordre infini. La description
géométrique est au § 9 (page 100) : un polygone inscrit dans une parabole.

\pagerange{96}{100}
\subsection*{La forme quadratique invariante et la conique circonscrite (pages 96 à 100)}

\textbf{§ 6.} Cherchons les fonctions polynomiales de degré $\leq 2$ sur $E$ ---
les formes quadratiques sur l'enveloppe vectorielle $\hat E$ --- invariantes
par $G = \mathrm{Aut}(\Pi_\infty)$, c'est-à-dire par $\sigma_0$ et $\sigma_1$.
À une constante près, l'invariance par $\sigma_1$ donne $f = a(x^2 + y^2) +
bxy + c(x + y)$, et l'invariance par $\sigma_0$ force $b = -\alpha a$, $c =
-a$ :
\[
  f = a\, f_\alpha, \qquad
  \boxed{\ f_\alpha(x, y) = (x^2 - x) + (y^2 - y) - \alpha\, xy\ } .
\]
Les fonctions $1$, $f_\alpha$ forment une base des fonctions invariantes de
degré $\leq 2$, et $f_\alpha$ est caractérisée par : a) $f_\alpha$ est nulle
sur $S$ ; b) sa partie homogène
\[
  f_0(X, Y) = X^2 + Y^2 - \alpha XY
\]
vaut $1$ sur les côtés, $f_0(s_i - s_{i-1}) = 1$. C'est la normalisation
« par les côtés » annoncée page 68. La \emph{conique circonscrite canonique}
$C(f_\alpha) : f_\alpha = 0$ ne dépend pas de la normalisation ; si le
polygone a au moins cinq sommets (en particulier pour les $n$-gones, $n \geq
5$, et les $\infty$-polygones), c'est la seule conique
circonscrite\footnote{La marge dit « si $\alpha \neq 0, 1$ (i.e.\ cas des
polygones réguliers à $n \geq 5$ côtés) ». Or $\alpha = 0$ est le carré et
$\alpha = -1$ le triangle (table de la page 120) ; $\alpha = 1$ est
l'hexagone. Il faut lire $\alpha \neq 0, -1$, ou mieux : au moins cinq
sommets distincts.}.

\textbf{§ 7. Discriminant de $f_0$.} Le déterminant de la matrice de $f_0$
(dans la convention $2 f_0$) est $\delta(f_0) = 4 - \alpha^2 = (2 - \alpha)(2
+ \alpha)$, nul si et seulement si $\alpha = \pm 2$ : $f_0$ est non dégénérée
exactement dans le cas ordinaire. Sur une extension contenant $\zeta$,
\[
  f_0(X, Y) = (Y - \zeta X)(Y - \zeta^{-1} X) ,
\]
et les deux directions isotropes $Y = \zeta X$, $Y = \zeta^{-1} X$ sont les
droites propres de $u_V$ --- ce qui traduit que $u_V$ est une
\emph{rotation} pour $f_0$. Pour $\alpha = \pm 2$ les deux directions se
confondent : $X + Y = 0$ pour $\alpha = -2$, la direction déjà privilégiée au
cas II ; $X = Y$ pour $\alpha = 2$. C'est la « métrique dégénérée » de la
page 52.

\textbf{§ 8. Discriminant de $f$.} La matrice de $2 f_\alpha$ en coordonnées
homogènes est
\[
  \begin{pmatrix} 2 & -\alpha & -1 \\ -\alpha & 2 & -1 \\ -1 & -1 & 0
  \end{pmatrix}, \qquad \det = -2(\alpha + 2) ,
\]
de sorte que, en caractéristique $\neq 2$, $C(f_\alpha)$ est non singulière si
et seulement si $\alpha \neq -2$\footnote{La page écrit $\delta = 2\delta'$,
$\delta' = -(\alpha + 2)$. Le calcul par la matrice symétrique suppose $2$
inversible ; la page ne traite pas à part la caractéristique $2$, où $-2 = 0$
et le cas $\alpha = -2$ est le carré.}. En particulier pour $\alpha = 2$, où
$f_0$ est dégénérée : $C$ est alors une \emph{parabole}. Pour $\alpha = -2$,
\[
  f_{-2}(x, y) = (x + y)^2 - (x + y) = (x + y)(x + y - 1) ,
\]
produit de deux formes affines de même partie homogène : $C$ est la réunion
des deux droites parallèles distinctes $x + y = 0$ et $x + y = 1$ qui portent
les sommets pairs et impairs (cas II) ; en coordonnées centrées, $X + Y = \mp
\frac12$.

\textbf{§ 9. Retour sur $\alpha = 2$} (caractéristique $\neq 2$). Avec $v = x
- y$, $w = x + y$,
\[
  f_2(x, y) = (x - y)^2 - (x + y) = v^2 - w ,
\]
et la rotation $u(x, y) = (1 - y + 2x, x)$ devient
\[
  u(v, w) = (v + 1,\ 1 + 2v + w), \qquad u(v, v^2) = (v + 1, (v + 1)^2) .
\]
Si l'on paramètre la parabole $C : w = v^2$ par $v$, $u|_C$ est la
\emph{translation} $v \mapsto v + 1$, et les sommets sont les points $(m,
m^2)$, $m \in \mathbf{Z}$, ou $m \in \mathbf{F}_p$ en caractéristique $p$ :
le $p$-gone régulier de caractéristique $p$ est formé des $p$ points
$\mathbf{F}_p$-rationnels d'une parabole, et sa rotation est une translation
le long de celle-ci\footnote{La page 100 s'arrête sur la formule et le point
d'exclamation « $v \mapsto v + 1$ ! » ; les deux dernières phrases (les
sommets $(m, m^2)$, le $p$-gone de $\mathbf{F}_p$) sont la conséquence immédiate,
écrite par nous, et rejoignent le cas III de la page 95.}.

\pagerange{101}{107}
\section*{IV. La fidélité (pages 101 à 107)}

Deux versions se suivent ; la seconde porte en marge « (ancienne version) »,
et la première, de la page 101, paraît donc la plus récente. Dans les deux,
$\Pi = (S, A, R)$ est un polygone combinatoire, $\rho = (\rho_S, \rho_A)$ une
réalisation géométrique régulière \emph{propre} (non impropre) de $\Pi$ dans un
plan $E$ sur $k$, et l'on note $S$, $A$ les images\footnote{Les pages
soulignent les lettres des objets combinatoires ($\underline{S}$,
$\underline{\Pi}$\ldots) et laissent nues leurs images, avec des
irrégularités ; on dit ici « combinatoire » ou « image » en toutes lettres.}.

\pagerange{101}{103}
\subsection*{Première version (pages 101 à 103)}

\textbf{Proposition.} L'image $\Pi_0 = (\rho_S(S), \rho_A(A))$ est un polygone
géométrique régulier de $E$ ; son groupe d'automorphismes $G_0$ est l'image de
$\mathrm{Aut}(\Pi)$ dans $\mathrm{Aff}(E)$, et $\Pi \to \Pi_0$ est un
morphisme de polygones combinatoires. En détail :
\begin{enumerate}
  \item[a)] Toute arête image $a = \mathrm{dr}(s, s')$ contient exactement deux
    sommets images. Les deux extrémités ont des images distinctes ; et un
    troisième sommet sur $a$ mettrait trois points de la conique circonscrite
    canonique $C$ sur une droite, ce qu'elle interdit, sauf si $C$ est
    singulière, $\alpha = -2$, cas où $C$ est la réunion de deux droites
    parallèles portant les sommets pairs et les impairs, et où
    $\mathrm{dr}(s_i, s_{i+1}) \cap C = \{s_i, s_{i+1}\}$\footnote{La page 101
    ne dit que « $a \cap C$ ne peut contenir que 2 points ($C$ conique
    circonscrite canonique) » ; le cas singulier est traité à la page 105, et
    on le rapatrie ici. Comme à la page 99, la caractéristique $2$ n'est pas
    examinée à part ; le carré de caractéristique $2$, où $f_0 = (x + y)(x +
    y + 1)$, relève du même argument que le cas $\alpha = -2$.}.
  \item[b)] Tout sommet image est sur exactement deux arêtes images : si
    $\mathrm{dr}(s_0, s_1)$, image d'une arête combinatoire, passe par $s$,
    alors $s \in \{s_0, s_1\}$ par a) ; si $s = s_0$, l'arête est l'image de
    l'une des deux arêtes combinatoires issues du sommet combinatoire au-dessus
    de $s$, à une rotation $u^n$ près qui fixe $s_0$, donc est l'identité.
  \item[c)] $S$ engendre $E$.
  \item[d)] $G = \mathrm{Aut}(E, S, A)$ est transitif sur l'ensemble $R$ des
    repères de $(S, A)$ : $R$ est l'image de l'ensemble des repères
    combinatoires, sur lequel $\mathrm{Aut}(\Pi)$ est transitif ; et comme
    $\mathrm{Aut}(E, S, A)$ y opère simplement transitivement, l'image de
    $\mathrm{Aut}(\Pi)$ est $\mathrm{Aut}(E, S, A)$ tout entier.
  \item[e)] $\Pi \to \Pi_0$ est un morphisme : l'application des repères est
    compatible à $\sigma_0$, $\sigma_1$.
\end{enumerate}

\pagerange{104}{107}
\subsection*{Seconde version, dite « ancienne » (pages 104 à 107)}

\textbf{Proposition.} Pour une réalisation régulière propre $\varphi$ de $\Pi$
sur $k$, les conditions suivantes sont équivalentes :
\begin{enumerate}
  \item[a)] $\varphi$ est fidèle (l'incidence des images entraîne
    l'incidence) ;
  \item[b)] $\varphi_S : S \to E$ est injective ; b$'$) $\varphi_A : A \to
    \mathrm{Dr}(E)$ est injective ;
  \item[c)] $\mathrm{Aut}(\Pi) \to \mathrm{Aff}(E)$ est injectif ;
  \item[d)] $\mathrm{Aut}^{+}(\Pi) \to \mathrm{Aff}(E)$ est injectif.
\end{enumerate}
\emph{Démonstration.} a) $\Rightarrow$ b), b$'$) : deux sommets distincts sont
séparés par une arête incidente à l'un et non à l'autre, puisque
$\operatorname{card} S \geq 3$. b) $\Rightarrow$ c) et b$'$) $\Rightarrow$ c) :
$\mathrm{Aut}(\Pi)$ s'injecte dans $\mathfrak{S}_S$ et dans $\mathfrak{S}_A$.
c) $\Rightarrow$ d) est trivial. d) $\Rightarrow$ b) : si $s_i = s_{i+d}$ avec
$0 < d < n$, alors $u^d s_0 = s_0$, donc $u^d$ fixe tous les $s_j$, qui
engendrent $E$, et $u^d = \mathrm{id}_E$, alors que $u^d \neq 1$ dans
$\mathrm{Aut}^{+}(\Pi)$. Enfin b), c), d) $\Rightarrow$ a) : si $s \in a =
\mathrm{dr}(s_0, s_1)$ avec $s \neq s_0, s_1$, la droite $a$ couperait $C$ en
trois points, donc $C$ serait singulière, $\alpha = -2$, et l'on a vu que
$\mathrm{dr}(s_0, s_1) \cap C = \{s_0, s_1\}$ : absurde.

\textbf{Corollaire} (page 105). Si $v \in \mathrm{Aut}^{+}(\Pi)$ et $s \in
S$, alors $v s = s$ si et seulement si $\varphi(v) = \mathrm{id}_E$ --- la
proposition 4 (a) de la page 22.

\textbf{Corollaire} (page 106). Soit $\varphi$ une réalisation quelconque d'un
polygone combinatoire $\Pi$ d'ordre $N$, $3 \leq N \leq \infty$, et $u$ l'image
d'un générateur de $\mathrm{Aut}^{+}(\Pi)$. Les entiers suivants sont égaux :
$\operatorname{card} \varphi(S)$, $\operatorname{card} \varphi(A)$,
$\operatorname{card} \varphi(\mathrm{Aut}^{+}(\Pi))$, l'ordre de $u$. Cet entier
$n$ est un diviseur de $N$ (condition vide si $N = \infty$), $n \geq 3$, et
c'est le plus petit diviseur $n$ de $N$ tel que $\varphi$ se factorise par le
quotient $\Pi_{(n)}$ de $\Pi$ d'ordre $n$, et l'unique tel que la réalisation
obtenue de $\Pi_{(n)}$ soit fidèle\footnote{La page avait d'abord écrit « le
plus grand », corrigé en « le plus petit ». C'est, sous forme intrinsèque, la
proposition 4 (c) de la page 23.}.

Mais il manque, dit la page, une théorie des morphismes (étales) de polygones
combinatoires qui permette de parler des quotients $\Pi_{(n)}$ de $\Pi$ pour
tout diviseur $n$ de l'ordre $N$. Elle n'est pas écrite\footnote{Un mot
illisible précède « comb. » ; le reste de la page 107 est blanc.}.

\pagerange{108}{120}
\section*{V. « Polynômes cyclotomiques » (pages 108 à 120)}

\pagerange{109}{113}
\subsection*{Les polynômes $F_n$ et le critère $u_\alpha^n = 1$ (pages 109 à 113)}

Soit $\mathbf{Z}[U] \subset \mathbf{Z}[T, T^{-1}]$, $U = T + T^{-1}$. On
définit, pour $n \geq 1$,
\[
  F_{2m}(U) = (1 + T^2 + \cdots + T^{2(m-1)})\, T^{-(m-1)}, \qquad
  F_{2m+1}(U) = (1 + T + \cdots + T^{2m})\, T^{-m} ,
\]
soit $F_{2m} = \dfrac{T^{2m} - 1}{T^2 - 1}\, T^{-(m-1)}$ et $F_{2m+1} =
\dfrac{T^{2m+1} - 1}{T - 1}\, T^{-m}$, polynômes unitaires en $U$ de degré
$[\frac{n-1}2]$ ; on pose $F_0 = 0$ et $F_{-n} = -F_n$. Les racines de $F_n$
sont les $\zeta + \zeta^{-1}$, $\zeta^n = 1$, $\zeta \neq \pm 1$, chaque paire
$\{\zeta, \zeta^{-1}\}$ une fois\footnote{Pour $n$ pair, $F_n(2\cos\theta) =
\sin(n\theta/2)/\sin\theta$, polynôme de Tchebychev de deuxième espèce en
$\cos\theta$ ; pour $n$ impair, $\sin(n\theta/2)/\sin(\theta/2)$, dit parfois
de « quatrième espèce ». Ces noms sont de nous. Le signe de « $F_n = -F_{-n}$ »
est peu net sur la page ; la formule
$F_n = (T^{n/2} - T^{-n/2})/(T^{\epsilon/2} - T^{-\epsilon/2})$,
$\epsilon \in \{1, 2\}$ selon la parité, le confirme.}. On a
\[
  F_{2m}(2) = m, \quad F_{2m}(-2) = (-1)^{m-1} m, \quad F_{2m+1}(2) = 2m + 1,
  \quad F_{2m+1}(-2) = (-1)^m ,
\]
et les premiers sont $F_1 = F_2 = 1$, $F_3 = U + 1$, $F_4 = U$, $F_5 = U^2 + U
- 1$, $F_6 = U^2 - 1$. Posons $X_n = F_n F_{n+1}$ (de degré $n - 1$).

\textbf{Proposition} (page 109). Pour tout $n \in \mathbf{Z}$,
\[
  X_n - U X_{n-1} + X_{n-2} = 1 , \qquad \text{c'est-à-dire} \qquad F_n
  F_{n+1} - U F_{n-1} F_n + F_{n-2} F_{n-1} = 1 .
\]
\emph{Démonstration.} Pour $n = 1, 2$ on vérifie directement ($1 - 0 + 0$,
$(U + 1) - U + 0$). Pour $n \geq 3$, on multiplie par $(T - 1)(T^2 - 1)
T^{n-1}$, en utilisant $F_n F_{n+1} (T - 1)(T^2 - 1) = (T^n - 1)(T^{n+1} - 1)
T^{-(n-1)}$ ; les deux membres deviennent $T^{n+2} - T^{n+1} - T^n + T^{n-1} =
(T - 1)(T^2 - 1) T^{n-1}$. Les $n$ négatifs s'en déduisent\footnote{La page
introduit d'abord des $X'_n = F_n F_{n+1}$, $Y'_n = F_n F_{n-1}$, puis des
$X_n$ définis par la récurrence $X_n = U X_{n-1} - X_{n-2} + 1$, et conclut
par récurrence $X_n = X'_n$ ; les primes manquent par endroits. On a fusionné
les deux notations. Une seconde démonstration (page 111) compare les racines :
$T^{n-1} X_n(U)$ est unitaire de degré $2(n-1)$ en $T$, de racines les $\xi
\neq 1$ avec $\xi^n = 1$ ou $\xi^{n+1} = 1$, au nombre de $2n - 2$ et
distinctes. Les deux lignes de la page 110, au pied d'un tapuscrit étranger,
sont $1 + T^2 + T^4 = T^2 (U^2 - 1)$ (le facteur $T^2$ manque sur la page) et
la formule de $F_{2m+2}$.}.

La proposition dit exactement que $F_{n-1}$ et $F_{n+1}$ engendrent l'idéal
unité :
\[
  A_n F_{n-1} + B_n F_{n+1} = 1, \qquad B_n = F_n, \quad A_n = F_{n-2} - U F_n
\]
\footnote{La page 113 écrit $A_n = F_{n-2} - U F_{n-1}$, l'indice du second
$F$ étant peu net ; le couple ne vérifie pas la relation ($n = 4$ : $(1 - U -
U^2)(U + 1) + U(U^2 + U - 1) = 1 - U - U^2$). Celui qu'on donne se lit sur
la proposition. La page 109 énonce $(F_n, F_{n+2}) = (1)$.}.

\textbf{Les sommets.} Dans les coordonnées universelles $(s_0, s_1, s_{-1})$
des pages 81 à 100, avec $u_\alpha(x, y) = (1 - y + \alpha x, x)$, la
récurrence $x_{n+1} = \alpha x_n - x_{n-1} + 1$, $y_{n+1} = x_n$, et les
valeurs initiales $s_0 = (0, 0)$, $s_1 = (1, 0)$, donnent
\[
  s_n = u_\alpha^n(s_0) = \bigl(F_n F_{n+1}(\alpha),\ F_{n-1} F_n(\alpha)\bigr)
  = \bigl(X_n(\alpha), X_{n-1}(\alpha)\bigr), \qquad n \in \mathbf{Z} .
\]
C'est la réponse à la « question » de la page 11 : le plan et les sommets
écrits en fonction de $\alpha$ seul, sur n'importe quel anneau. L'application
$u^n$ envoie $(s_0, s_1, s_{-1})$ sur $(s_n, s_{n+1}, s_{n-1})$, d'où sa
matrice
\[
  u_\alpha^n = \begin{pmatrix}
    F_{n+1}(F_{n+2} - F_n) & F_n (F_{n-1} - F_{n+1}) & F_n F_{n+1} \\
    F_n (F_{n+1} - F_{n-1}) & F_{n-1}(F_{n-2} - F_n) & F_{n-1} F_n \\
    0 & 0 & 1
  \end{pmatrix}(\alpha)
\]
\footnote{La page donne deux versions ; la première, barrée, a la troisième
colonne $(F_n F_{n+1}, F_{n-1} F_n)$, la seconde $(F_{n+1} F_{n+1}, F_{n-1}
F_{n+1})$, indices surchargés et douteux. La troisième colonne est $s_n$, et
c'est la première version qui est juste.}.

\textbf{Théorème} (page 112). Soit $k$ un anneau, $\alpha \in k$, $u_\alpha$
l'automorphisme affine $(x, y) \mapsto (1 - y + \alpha x, x)$ de $k^2$, $s_n =
u_\alpha^n(0, 0)$, et $n \in \mathbf{Z}$. Les conditions suivantes sont
équivalentes :
\begin{enumerate}
  \item[a)] $u_\alpha^n = \mathrm{id}$ ;
  \item[b)] $s_n = s_0$ ;
  \item[c)] $u_\alpha^n s_i = s_i$, pour un $i \in \mathbf{Z}$ fixé ;
  \item[d)] $F_n(\alpha) = 0$.
\end{enumerate}
\emph{Démonstration.} a) $\Rightarrow$ b) est trivial ; b) entraîne $u^n s_i =
s_i$ pour tout $i$ (multiplier par $u^i$), en particulier pour $i = \pm 1$, et
$s_0, s_1, s_{-1}$ forment un repère affine, d'où a) ; b) $\Leftrightarrow$ c)
en multipliant par $u^{\pm i}$. d) $\Rightarrow$ b) par la formule des sommets.
Réciproquement, si $F_n F_{n+1}(\alpha) = F_{n-1} F_n(\alpha) = 0$, la
relation de Bézout donne $F_n(\alpha) = F_n(\alpha)(A_n F_{n-1} + B_n
F_{n+1})(\alpha) = 0$\footnote{La page conclut « la relation $F_n(\alpha) =
1$, cqfd » : lire $= 0$. Elle lettre les conditions a), b), b$_i$), d), et
conclut « a $\iff$ b $\iff$ c », où c) est b$_i$).}.

Pour $\alpha = 2$ en caractéristique $p$ impaire, $F_p(2) = p = 0$ et
$u^p = 1$ ; pour $\alpha = -2$, $F_n(-2) = 0$ demande $n = 2m$ pair avec $p
\mid m$, d'où l'ordre $2p$ : on retrouve les pages 94 et 95. Par la
factorisation $F_n = \prod_{d \mid n,\ d \geq 3} \Psi_d$, le théorème dit, sur
un corps de caractéristique $0$, que $u_\alpha^n = 1$ si et seulement si
$\alpha$ est l'invariant d'un $d$-gone régulier pour un diviseur $d \geq 3$ de
$n$.

\pagerange{114}{116}
\subsection*{Le polynôme semi-cyclotomique et les anneaux modulaires (pages 114 à 116)}

Soit $\Phi_n(T) = \prod (T - \xi)$, $\xi$ parcourant les racines primitives
$n$-ièmes de $1$ dans $\bar{\mathbf{Q}}$ : polynôme unitaire de
$\mathbf{Z}[T]$, de degré $\varphi(n) = n \prod_{p \mid n}(1 - \frac1p)$, avec
$T^n - 1 = \prod_{d \mid n} \Phi_d$. Sur un corps de caractéristique $p \nmid
n$, les racines de $\Phi_n$ sont les racines primitives $n$-ièmes de $1$.

Si une réalisation régulière de $\Pi_\infty$ sur $k$ ($p \nmid n$) est donnée
par $\alpha = \xi + \xi^{-1}$, elle se factorise en une réalisation
\emph{fidèle} de $\Pi_n$, $n \geq 3$, si et seulement si $\Phi_n(\xi) = 0$,
c'est-à-dire $\Psi_n(\alpha) = 0$, où le \emph{polynôme semi-cyclotomique}
$\Psi_n \in \mathbf{Z}[U]$ est défini par
\[
  \Psi_n(T + T^{-1}) = T^{-\varphi(n)/2}\, \Phi_n(T) \qquad (n \geq 3) :
\]
c'est le polynôme minimal de $2\cos(2\pi/n)$\footnote{Nom moderne de nous ;
ce polynôme est étudié au moins depuis D. H. Lehmer (1933), sous le nom de
polynôme minimal de $\cos(2\pi/n)$ ou de $2\cos(2\pi/n)$ (cité de mémoire).
La page écrit $\mathbf{Z}(U)$ pour $\mathbf{Z}[U]$.}.

\textbf{Les anneaux modulaires.} Si $n \geq 3$ n'est ni premier ni double d'un
premier ($n = 8, 9, 12, 15, 16, \ldots$), l'anneau « modulaire » des
réalisations fidèles de $\Pi_n$ est
\[
  k'_n = \mathbf{Z}[\tfrac1n][U]/\Psi_n(U) ,
\]
l'invariant universel $\alpha'_n$ étant l'image de $U$ ; il faut inverser $n$,
parce qu'en caractéristique $p \mid n$ une racine de $\Psi_n$ définit un
polygone d'un autre ordre (corollaire de la page 118). Pour les autres $n$, on
peut faire mieux :
\begin{itemize}
  \item \emph{$n = p$ premier impair.} Toute réalisation régulière de $\Pi_p$
    sur un corps est fidèle (page 27, corrigée) ; $\Phi_p = 1 + T + \cdots +
    T^{p-1}$, $\Psi_p = F_p$, et $k'_p = \mathbf{Z}[U]/F_p(U)$, sans rien
    inverser : en caractéristique $p$, la seule racine est $2$, le $p$-gone
    parabolique.
  \item \emph{$n = 2p$, $p$ premier impair.} $\Phi_{2p}(T) = \Phi_p(-T)$,
    $\Psi_{2p}(U) = (-1)^{(p-1)/2}\,\Psi_p(-U)$. En caractéristique $p$, les
    racines de $\Phi_{2p}$ sont toutes égales à $-1$, celles de $\Psi_{2p}$ à
    $-2$, invariant du $2p$-gone en bande ; en caractéristique $2$, en
    revanche, elles définissent des $p$-gones. Donc en toute caractéristique
    $\neq 2$, $\Psi_{2p}(\alpha) = 0$ est la condition nécessaire et suffisante
    pour que $\alpha$ définisse un $2p$-gone régulier, et $k'_{2p} =
    \mathbf{Z}[\frac12][U]/\Psi_{2p}(U)$\footnote{La page 116 écrit que $-2$
    « est bien l'invariant du pol.\ régulier à $2$ côtés » : il s'agit du
    $2p$-gone. « Nécessaire et suffisante » est une lecture de l'abréviation
    « n.\ et s. ». Le passage barré de la page 115 (racines de $\Phi_p$ en
    caractéristique $p$) n'est pas repris.}.
  \item \emph{$n = 4$.} $\Phi_4 = T^2 + 1$, $\Psi_4 = U = F_4$ ; en
    caractéristique $2$, le seul zéro $0 = -2$ définit bien le carré. Donc
    $k'_4 = \mathbf{Z}[U]/(U) \simeq \mathbf{Z}$ : sur un anneau donné il y a
    une et une seule réalisation régulière de $\Pi_4$, à isomorphisme unique
    près, fidèle quand l'anneau est un corps.
\end{itemize}
C'est la réponse partielle à la question b) de la page 47 : le carré est
défini sur $\mathbf{Z}$ ; pour $n = p$, l'anneau $\mathbf{Z}[U]/F_p$ n'a pas à
inverser $p$.

\pagerange{117}{120}
\subsection*{Réduction modulo $p$, ordre en caractéristique $p$, table (pages 117 à 120)}

\textbf{Proposition} (page 117). Soient $p$ premier, $n = p^r n'$ avec $(p,
n') = 1$. Alors
\[
  \Phi_n(T) \equiv \Phi_{n'}(T)^{\varphi(p^r)} \pmod p .
\]
\emph{Démonstration}, par récurrence sur $r + n'$. Pour tout diviseur strict
$\delta = p^s d$ de $n$, l'hypothèse de récurrence donne $\Phi_\delta \equiv
\Phi_d^{\varphi(p^s)}$. Multipliant ces relations entre elles et par
$\Phi_n$,
\[
  \Phi_{n'}^{\varphi(p^r)} \prod_{\delta \mid n} \Phi_\delta \equiv
  \prod_{s = 0}^{r} \Bigl(\prod_{d \mid n'} \Phi_d\Bigr)^{\varphi(p^s)}
  \Phi_n , \quad\text{soit}\quad
  \Phi_{n'}^{\varphi(p^r)} (T^n - 1) \equiv (T^{n'} - 1)^{p^r}\, \Phi_n ,
\]
puisque $\sum_{s=0}^{r} \varphi(p^s) = p^r$. Or $T^n - 1 = (T^{n'})^{p^r} -
1 \equiv (T^{n'} - 1)^{p^r}$, et l'on simplifie dans l'anneau intègre
$\mathbf{F}_p[T]$\footnote{La formule (2) de la page a été corrigée en place et
se réduit à $\Phi_n \equiv \Phi_n$ ; on en donne le rôle. Les indices des deux
$\Phi$ de la conclusion sont mal formés ; l'énoncé les fixe.}.

\textbf{Corollaire} (pages 118 et 119). Soit $k$ un corps de caractéristique
$p > 0$, $n \geq 1$, $\alpha \in k$ avec $\Psi_n(\alpha) = 0$, $n = p^r n'$,
$(p, n') = 1$. Le polygone régulier d'invariant $\alpha$ est d'ordre
\[
  n' \text{ si } n' \geq 3 ; \qquad 2p \text{ si } n' = 2 ; \qquad p \text{ si
  } n' = 1 \text{ et } p \neq 2 ; \qquad 4 \text{ si } n' = 1 \text{ et } p = 2 .
\]
\emph{Démonstration.} Écrivons $\alpha = \xi + \xi^{-1}$, $\xi$ dans une
extension. $\Psi_n(\alpha) = 0$ équivaut à $\Phi_n(\xi) = 0$, donc, par la
proposition, à $\Phi_{n'}(\xi) = 0$ : $\xi$ est racine primitive $n'$-ième de
$1$. Si $n' \geq 3$, $\alpha \neq \pm 2$ et l'ordre est celui de $\xi$ ; si
$n' = 2$, $\alpha = -2$ et l'ordre est $2p$ (cas II) ; si $n' = 1$, $\alpha =
2$ et l'ordre est $p$, ou $4$ si $p = 2$ (cas III)\footnote{La page dit
« d'ordre $p$ si $n' = 1$ », sans exception : pour $p = 2$, $\alpha = 2 = 0$
est le carré, d'ordre $4$, comme la page 95 l'a établi. La page avait d'abord
écrit $n \geq 3$, corrigé en $n \geq 1$.}.

\textbf{La table de la page 120.} $\varphi(n)$ est pair pour $n \geq 3$, et
$\varphi(p^r) < 6$ seulement pour $p^r = 2, 3, 4, 5, 8$. D'où les petits
degrés :
\begin{itemize}
  \item $\varphi(n) \leq 2$, $\alpha$ \emph{entier} : $n = 1, 2, 4, 3, 6$,
    $\alpha = 2, -2, 0, -1, 1$ ;
  \item $\varphi(n) = 4$, $\alpha$ \emph{quadratique} : $n = 8$, $\alpha^2 = 2$ ;
    $n = 12$, $\alpha^2 = 3$ ; $n = 5$, $\alpha^2 + \alpha - 1 = 0$ ; $n = 10$,
    $\alpha^2 - \alpha - 1 = 0$ ;
  \item $\varphi(n) = 6$, $\alpha$ \emph{cubique} : $n = 9$, $\alpha^3 - 3\alpha
    + 1 = 0$ ; $n = 7$, $\alpha^3 + \alpha^2 - 2\alpha - 1 = 0$ ; $n = 18$,
    $\alpha^3 - 3\alpha - 1 = 0$ ; $n = 14$, $\alpha^3 - \alpha^2 - 2\alpha + 1 =
    0$.
\end{itemize}
Les seuls polygones réguliers définis sur $\mathbf{Q}$ sont donc le triangle,
le carré et l'hexagone. Les cubiques sont, dit une note lue avec doute,
des « exemples d'équations du 3e degré non résolubles » --- trois
racines réelles, irréductibles sur $\mathbf{Q}$ : le \emph{casus
irreducibilis}, où les radicaux réels ne suffisent pas\footnote{L'expression
latine est de nous ; la note de la page est à demi illisible. La marge
justifie la liste $\varphi(n) = 6$ : $p \neq 2$, et $r = 1$, $p = 7$ ou $r =
2$, $p = 3$.}.

\pagerange{121}{129}
\section*{VI. Feuilles détachées (pages 121 à 129)}

\pagerange{121}{121}
\subsection*{Un calcul tensoriel barré (page 121)}

La page, barrée de deux longs traits, commence au milieu d'un calcul : avec la
transposée ${}^t u$ et l'identité $({}^t u\, x') \otimes x = x' \otimes u x$
dans $V^{*} \otimes V$, il cherche la forme des matrices $u$ qui satisfont une
relation sur l'élément $(-2e'_1 + e'_2) \otimes e_1$, et trouve une matrice
triangulaire de diagonale $\lambda$ ; tête-bêche, un polynôme caractéristique
$(\lambda - 1)^{\nu - 1}(\lambda + 1)$, celui d'une réflexion en dimension
$\nu$. Le but du calcul ne se lit pas, et l'on n'en tire rien\footnote{Deux
croquis de trapèze, et en marge des fragments ($x + \langle x, x' \rangle x$,
$-x$), suggèrent des réflexions $x \mapsto x - \langle x, x' \rangle a$ ; c'est
une conjecture qu'on n'appuie sur rien.}.

\pagerange{122}{124}
\subsection*{Tables et réductions modulo $p$ (pages 122 à 124)}

La page 122 reprend la matière des pages 114 à 120 avec d'autres noms :
$\Phi_n$ y est le polynôme minimal sur $\mathbf{Q}$ d'une racine primitive
$n$-ième de $1$, et ce qui est ici $\Psi_n$ le polynôme minimal de $\zeta +
\zeta^{-1}$, avec $\Phi_n(T) = T^{\varphi(n)/2}\, \Psi_n(T + T^{-1})$ pour $n
\neq 1, 2$, $\Psi_1 = U - 2$, $\Psi_2 = U + 2$\footnote{La page appelle ce
polynôme $F_n$, alors que le $F_n$ des pages 109 à 113 est autre ($F_6 = U^2 -
1$ à la page 109, « $F_6 = U - 1$ » ici). On écrit $\Psi_n$ ; dans la table,
la variable est une capitale barrée, lue $U$.}. La table :
\[
\begin{array}{lll}
\Psi_1 = U - 2, & \Psi_2 = U + 2, & \Psi_3 = U + 1, \\
\Psi_4 = U, & \Psi_5 = U^2 + U - 1, & \Psi_6 = U - 1, \\
\Psi_7 = U^3 + U^2 - 2U - 1, & \Psi_8 = U^2 - 2, & \Psi_9 = U^3 - 3U + 1, \\
\Psi_{10} = U^2 - U - 1, & &
\end{array}
\]
en regard de $\Phi_1, \ldots, \Phi_{10}$ ; « donner aussi les cas $n = 12, 14,
18$ » --- ceux de la page 120. Pour $p = 2p' + 1$ premier impair,
\[
  \Psi_p = S_0 + S_1 + \cdots + S_{p'}, \qquad
  \Psi_{2p}(U) = (-1)^{p'}\, \Psi_p(-U) = (-1)^{p'} \sum_{j=0}^{p'} (-1)^j S_j(U),
\]
où $S_0 = 1$, $S_1 = U$, $S_2 = U^2 - 2$, $S_3 = U^3 - 3U$, \ldots, et
\[
  U^n = \sum_{0 \leq i \leq [n/2]} \binom{n}{i} S_{n-2i}(U), \qquad n \geq 1 ,
\]
d'où le calcul récurrent des $S_n$\footnote{Pour $n$ pair, le terme $i = n/2$
vaut $\binom{n}{n/2} S_0 = \binom{n}{n/2}$, ce qui est juste puisque $S_0 =
1$ et non $T^0 + T^0 = 2$ : la formule est correcte telle qu'elle est écrite.
Les $S_n$ sont, à la normalisation près, les polynômes de Tchebychev de
première espèce ($S_n(2\cos\theta) = 2\cos n\theta$), ou les polynômes de
Dickson $D_n(U, 1)$ (noms de nous). À la page 123, la lettre devant $(-T)$ est
lue $\Phi_p$ alors que la formule appelle $\Psi_p$.}.

\textbf{Réduction modulo $p$} (pages 123 et 124). Sur un corps algébriquement
clos de caractéristique $p$, avec $n = p^r n'$, $(p, n') = 1$, les racines de
$\Phi_n$ sont les racines primitives $n'$-ièmes de $1$, chacune de
multiplicité $\varphi(p^r)$ --- c'est la proposition de la page 117. Pour
$\Psi_n$ ($n \geq 3$) :
\begin{enumerate}
  \item[1)] si $n' \neq 1, 2$, $\Psi_n \equiv \Psi_{n'}^{\varphi(p^r)} \pmod
    p$ ;
  \item[2)] si $n' = 1$ ou $2$, $\Psi_n \equiv \Psi_{n'}^{\varphi(p^r)/2}$,
    c'est-à-dire $(U - 2)^{\varphi(p^r)/2}$ si $n = p^r$ et $(U +
    2)^{\varphi(p^r)/2}$ si $n = 2p^r$.
\end{enumerate}
En particulier, pour $p$ premier impair et $p' = (p - 1)/2$,
\[
  \Psi_p \equiv (U - 2)^{p'}, \qquad \Psi_{2p} \equiv (U + 2)^{p'} \pmod p
\]
\footnote{Variances : la page écrit la condition « $(n, n') = 1$ » pour $(p,
n') = 1$ ; en 2), « $p^r \neq 2$ i.e.\ $n \neq 4$ », alors que le cas exclu est
$n = 2$ (pour $n = 4$, $p = 2$, $\Psi_4 = U \equiv U - 2$ et la formule vaut) ;
elle écrit d'abord $F_{p^r} \equiv (T - 2)^{\varphi(p^r)}$, exposant que la
page 124 corrige en $\varphi(p^r)/2$ ; et, à la fin de la page 124, $F_p
\equiv (T - 1)^{p'}$, $F_{2p} \equiv (T + 1)^{p'}$, où la formule générale de
la même page donne $T - 2$ et $T + 2$.}. C'est la forme polynomiale du
corollaire de la page 118 : en caractéristique $p$, le $p$-gone et le
$2p$-gone n'ont plus qu'un invariant, $2$ et $-2$.

\pagerange{125}{127}
\subsection*{Les anneaux $\mathbf{Z}[\zeta]$, $\mathbf{Z}[\alpha]$ et la ramification en $p$ (pages 125 à 127)}

Pour $p$ premier impair, soient
\[
  A_p = \mathbf{Z}[T]/\Phi_p(T) \ni \zeta, \qquad B_p = \mathbf{Z}[U]/\Psi_p(U)
  \ni \alpha = \zeta + \zeta^{-1} ,
\]
anneaux \emph{normaux} --- l'anneau des entiers du corps cyclotomique
$\mathbf{Q}(\zeta_p)$ et celui de son sous-corps réel maximal, dirait-on
aujourd'hui\footnote{La normalité est affirmée, non démontrée ; c'est un
théorème classique. Les noms sont de nous. L'indice de $B$ est surchargé
(« $n(p)$ » sous un $p$).}. Ils sont étales sur $\mathbf{Z}$ en dehors de $p$,
les racines de $\Phi_p$, $\Psi_p$ étant distinctes en caractéristique $\ell
\neq p$, et ramifiés en $p$, où il n'y a plus qu'une racine. Dans $A_p$,
$(\zeta - 1)^{p-1} \equiv 0 \pmod p$ ; et la norme $N(\zeta - 1) =
\prod_{i}(\zeta_i - 1) = \Phi_p(1) = p$ (au signe près) donne
\[
  \boxed{\ (\zeta - 1)^{p-1} = p\, \varepsilon\ }, \qquad \varepsilon \text{
  unité} ,
\]
d'où $A/(\zeta - 1) = \mathbf{F}_p[T]/(T - 1)^{p-1}/(T - 1) = \mathbf{F}_p$ :
$(\zeta - 1)$ est un idéal maximal, et $p$ est \emph{totalement ramifié} dans
$A_p$\footnote{La page écrit $F_p(\zeta)$ et $F_p(1)$ là où le calcul porte sur
$\Phi_p$. « Totalement ramifié » est de nous ; la page dit « ramifiés en
$\ell = p$ (une seule racine en car.\ $p$) ».}. De même, pour $p^r \neq 2$ et
$\varphi' = \varphi(p^r)/2$, dans $B_{p^r} = \mathbf{Z}[U]/\Psi_{p^r}$,
\[
  (\alpha - 2)^{\varphi'} = p\, \varepsilon', \qquad N(\alpha - 2) = \pm
  \Psi_{p^r}(2) = \pm \Phi_{p^r}(1) = \pm p ,
\]
puisque $\prod_{1 \neq d \mid p^r} \Phi_d(1) = 1 + 1 + \cdots + 1 = p^r$
donne $\Phi_{p^r}(1) = p$ (page 127). La page vérifie au passage
$\Psi_7(2) = 8 + 4 - 4 - 1 = 7$, et s'interroge (« $F_p(2) = p$ ? ») --- ce
qui est le cas, $F_p(2) = p$ par la page 109. Enfin
\[
  (\zeta - 1)(\zeta^{-1} - 1) = 2 - \alpha ,
\]
donc $\alpha - 2$ est, à une unité près, le carré de $\zeta - 1$ : la place
au-dessus de $p$ dans $B$ se ramifie encore dans $A$, et la page demande si
$B[T]/(T^2 - \alpha T + 1) \to A$ est un isomorphisme --- ce que la page 128
établit\footnote{Les pages 126 et 127 sont des brouillons serrés ; la page 126
esquisse le cas général $n = n' p^r$ ($\mathbf{F}_p[T]/\Phi_n =
\mathbf{F}_p[T]/\Phi_{n'}^{\varphi(p^r)}$, la norme relative
$N_{A_n/A_{n'}}(\zeta - 1)$, $\zeta' = \zeta^{p^r}$) sans conclure ; on n'en
retient que ces formules.}.

\pagerange{128}{128}
\subsection*{Les anneaux $A_n$, $B_n$ et l'involution $T \mapsto T^{-1}$ (page 128)}

Soient $A_n = \mathbf{Z}[T]/\Phi_n(T)$ et $B_n = \mathbf{Z}[U]/\Psi_n(U)$. Pour
$n = n' n''$ avec $(n', n'') = 1$, $(\mathbf{Z}/n)^{*} \simeq
(\mathbf{Z}/n')^{*} \times (\mathbf{Z}/n'')^{*}$ et $A_{n'} \otimes_{\mathbf{Z}}
A_{n''} \xrightarrow{\sim} A_n$ ; en termes de schémas, le schéma des racines
primitives $\mu_n^{*} = \mathrm{Spec}\, A_n \subset \mu_n$ vérifie
$\mu^{*}_{n'n''} \simeq \mu^{*}_{n'} \times \mu^{*}_{n''}$, et $\mu_n^{*}/\{\pm
1\} = \mathrm{Spec}\, B_n$. La page en tire la normalité de $A_n$ « par
hypothèse de récurrence »\footnote{L'argument est esquissé. Un produit
tensoriel d'anneaux normaux n'est pas normal en général ; il l'est ici parce
que $A_{n'}$ est étale au-dessus des premiers où $A_{n''}$ se ramifie, et
inversement --- ce que la page ne dit pas, et qui est peut-être l'objet de la
page 129.}.

L'involution $\sigma : T \mapsto T^{-1}$ de $A = \mathbf{Z}[T, T^{-1}]$ a pour
invariants $B = \mathbf{Z}[U]$ ; géométriquement, le quotient de
$\mathbf{G}_m$ par l'inversion est la droite affine $\mathbf{E}^1$, de
coordonnée $U = T + T^{-1}$ :
\begin{tikzcd}
A = \mathbf{Z}[T, T^{-1}] \arrow[r] & A_n = \mathbf{Z}[T]/\Phi_n \\
B = \mathbf{Z}[U] = A^{\sigma} \arrow[r] \arrow[u, hook] & B_n = \mathbf{Z}[U]/\Psi_n \arrow[u, hook]
\end{tikzcd}
Pour $n \neq 1, 2$, $A = B[S]/(S^2 - US + 1)$ est libre sur $B$ de base $1,
T$ : si $F(U) + T\,G(U) = 0$, alors, appliquant $\sigma$, $F(U) + T^{-1} G(U)
= 0$, d'où $(T - T^{-1}) G(U) = 0$ et $G = 0$. De même $A_n = B_n \oplus B_n
T$, avec $\sigma(F + GT) = (F + GU) - GT$ (car $T^{-1} = U - T$), et
$B_n = A_n^{\sigma}$ : $\sigma(F + GT) = F + GT$ force $2G = 0$, donc $G = 0$,
$B_n$ étant sans torsion\footnote{La page écrit la condition « $GU = 0$, $G =
-G$ $\iff G = 0$ ». L'application $B_n[T]/(T^2 - \alpha_n T + 1) \to A_n$,
$T \mapsto \xi_n$, répond à la question de la page 127. L'exposant de $A_n$
dans « $B_n = A_n^{\sigma}$ » est lu avec doute.}.

\pagerange{129}{129}
\subsection*{Revêtements ramifiés (page 129)}

La dernière page change de langage. Sur une base $S$, soient $X$ et $X'$ deux
revêtements normaux irréductibles, de groupes $G$ et $G'$, $X$ totalement
ramifié en un point $s$ et étale ailleurs, $X'$ totalement ramifié en $s' \neq
s$ et étale ailleurs. Le produit fibré $X \times_S X'$ est-il connexe ? Il
dessine le carré des groupes fondamentaux
\begin{tikzcd}[column sep=small, row sep=small]
& \pi_1(S \smallsetminus \{s, s'\}) \arrow[dl] \arrow[dr] & \\
\pi_1(S \smallsetminus s) \arrow[dr] & & \pi_1(S \smallsetminus s') \arrow[dl] \\
& \pi_1(S) &
\end{tikzcd}
et un modèle où ôter un point ajoute un générateur libre ($G * \mathbf{Z} *
\mathbf{Z} \to G * \mathbf{Z} \to G$), puis un sous-groupe d'inertie $I
\subset H_s \subset G \times G'$ se surjectant sur $G$ et sur $G'$, et un
revêtement $Z \subset X \times_S X'$ totalement ramifié au-dessus de $X_s$ et
de $X'_{s'}$\footnote{Un sommet du diagramme se lit « $X * \mathbf{Z}$ », où
l'on attend $G * \mathbf{Z}$ ; la dernière ligne écrit $X$ pour $X'$ ; le
reste (« $(\mathbf{Z}/4)^{*}$ d'ordre $2$ ») ne se raccorde à rien.}. Le
raisonnement reste en suspens, et le dossier s'arrête là.

C'est, selon toute vraisemblance, la forme géométrique de la question des
pages 127 et 128 : $\mathrm{Spec}\, A_{n'}$ et $\mathrm{Spec}\, A_{n''}$ sont
des revêtements de $\mathrm{Spec}\, \mathbf{Z}$ ramifiés en des ensembles
disjoints de premiers, et la connexité (avec la normalité) de leur produit
fibré, $\mathrm{Spec}(A_{n'} \otimes A_{n''})$, est ce qu'il faut pour que
$A_{n'} \otimes A_{n''} \simeq A_{n'n''}$ soit intègre et normal ; c'est le
mécanisme où intervient, en arithmétique, le fait que $\mathrm{Spec}\,
\mathbf{Z}$ n'a pas de revêtement étale non trivial (théorème de
Minkowski)\footnote{Ce rapprochement est entièrement de nous ; rien sur la page
ne nomme $\mathbf{Z}$ ni les anneaux $A_n$, et le lien n'est suggéré que par
le voisinage de la page 128.}.

\end{document}
